% This chapter has been modified on 6-4-05.
%\setcounter{chapter}{5}


\chapter{Nilai Harapan dan Varians}\label{chp 6}

\section[Nilai Harapan]{Nilai Harapan Peubah Acak Diskret}\label{sec 6.1}

Ketika sekumpulan besar bilangan dihimpun, seperti dalam suatu sensus, biasanya kita
tidak tertarik pada tiap-tiap bilangan, melainkan pada besaran deskriptif tertentu,
seperti rata-rata atau median. Secara umum, hal yang sama berlaku bagi distribusi
peluang suatu peubah acak bernilai numerik. Dalam bagian ini dan bagian berikutnya,
kita akan membahas dua besaran deskriptif semacam itu: \emx {nilai harapan} dan
\emx {varians.} Kedua besaran ini hanya berlaku bagi peubah acak bernilai numerik;
karena itu, dalam bagian-bagian ini kita menganggap bahwa semua peubah acak mempunyai
nilai numerik. Untuk memberikan pembenaran intuitif bagi definisi kita, tinjaulah
permainan berikut.

\subsection*{Nilai Rata-Rata}

Sebuah dadu dilempar. Jika muncul bilangan ganjil, kita menang sebesar bilangan
tersebut; jika muncul bilangan genap, kita kalah sebesar bilangan tersebut. Sebagai
contoh, jika muncul dua kita kalah 2, dan jika muncul tiga kita menang 3. Kita ingin
menentukan apakah permainan ini layak dimainkan. Mula-mula kita mencoba simulasi.
Program {\bf Die}\index{Die (program)} menjalankan simulasi ini.
\par
Program tersebut mencetak frekuensi dan frekuensi relatif munculnya setiap hasil.
Program itu juga menghitung rata-rata kemenangan. Kita telah menjalankan program dua
kali. Hasilnya ditampilkan dalam Tabel~\ref{table 6.1}.

\medskip
\begin{table}[h]\centering
\centering
\begin{tabular}{|c|c|c|c|c|} \hline
        & \multicolumn{2}{c|}{n = 100}    & \multicolumn{2}{c|}{n = 10000} \\ \cline{2-5}
Kemenangan & \hskip.1in Frekuensi \hskip.1in & Relatif  & \hskip.1in Frekuensi \hskip.1in & Relatif  \\  
        &                                 & Frekuensi &                                 & Frekuensi
\\ \cline{1-5} 1       & 17        & .17                 & 1681       &.1681 \\  -2      & 17       
& .17                 & 1678       &.1678 \\   3       & 16        & .16                 &
1626       &.1626 \\   -4      & 18        & .18                 & 1696       &.1696 \\ 
5       & 16        & .16                 & 1686       &.1686 \\ 
-6      & 16        & .16                 & 1633       &.1633 \\\hline
\end{tabular}
\caption{Frekuensi untuk permainan dadu.}
\label{table 6.1}
\end{table}

Pada pengoperasian pertama, kita memainkan permainan itu 100 kali. Dalam pengoperasian
ini, rata-rata keuntungan kita adalah $-.57$. Tampaknya permainan itu merugikan, dan
kita ingin mengetahui seberapa merugikannya. Untuk memperoleh gambaran yang lebih
baik, kita memainkan permainan itu 10{,}000 kali. Dalam kasus ini, rata-rata keuntungan
kita adalah $-.4949$.
\par
Kita perhatikan bahwa frekuensi relatif setiap satu dari enam hasil yang mungkin cukup
dekat dengan peluang 1/6 bagi hasil tersebut. Hal ini sesuai dengan penafsiran
frekuensi kita atas peluang. Hal ini juga menunjukkan bahwa untuk jumlah permainan
yang sangat besar, rata-rata keuntungan kita seharusnya
\begin{eqnarray*}
\mu & = & 1 \Bigl(\frac 16\Bigr) - 2\Bigl(\frac 16\Bigr) + 3 \Bigl(\frac 16\Bigr) - 4
\Bigl(\frac 16\Bigr) + 5 \Bigl(\frac 16\Bigr) - 6 \Bigl(\frac 16\Bigr) \\
    & = & \frac 96 - \frac {12}6 = -\frac 36 = -.5\ .
\end{eqnarray*} Nilai ini cukup sesuai dengan rata-rata keuntungan kita untuk 10{,}000
permainan.

Kita perhatikan bahwa nilai yang kita pilih sebagai rata-rata keuntungan diperoleh
dengan mengambil hasil-hasil yang mungkin, mengalikan masing-masing dengan peluangnya,
lalu menjumlahkan hasil perkalian itu. Hal ini menyarankan definisi berikut bagi hasil
yang diharapkan dari suatu percobaan.

\subsection*{Nilai Harapan}

\begin{definition}\label{def 6.1} Misalkan $X$ adalah peubah acak diskret bernilai
numerik dengan ruang sampel $\Omega$ dan fungsi distribusi $m(x)$. {\em Nilai
harapan}\index{nilai harapan} $E(X)$ didefinisikan oleh
$$ E(X) = \sum_{x \in \Omega} x m(x)\ ,
$$ asalkan jumlah ini konvergen mutlak. Nilai harapan sering kita sebut
\emx {rataan,}\index{rataan} dan untuk singkatnya $E(X)$ kita lambangkan dengan
$\mu$. Jika jumlah di atas tidak konvergen mutlak, kita katakan bahwa $X$ tidak
mempunyai nilai harapan.
\end{definition}

\begin{example}\label{exam 6.03} Misalkan suatu percobaan terdiri atas pelemparan
sebuah koin adil sebanyak tiga kali. Misalkan $X$ menyatakan banyaknya kepala yang
muncul. Nilai-nilai $X$ yang mungkin adalah $0, 1, 2$, dan $3$. Peluang yang
bersesuaian adalah $1/8, 3/8, 3/8,$ dan $1/8$. Dengan demikian, nilai harapan $X$
adalah
$$0\biggl(\frac 18\biggr) + 1\biggl(\frac 38\biggr) + 2\biggl(\frac 38\biggr) +
3\biggl(\frac 18\biggr) = \frac 32\ .$$ Nanti dalam bagian ini kita akan melihat
cara yang lebih cepat untuk menghitung nilai harapan tersebut, berdasarkan fakta bahwa
$X$ dapat ditulis sebagai jumlah peubah-peubah acak yang lebih sederhana.
\end{example}


\begin{example}\label{exam 6.05} Misalkan kita melempar sebuah koin adil hingga
kepala muncul untuk pertama kalinya, dan misalkan $X$ menyatakan banyaknya lemparan
yang dilakukan. Nilai-nilai $X$ yang mungkin adalah $1, 2, \ldots$, dan fungsi
distribusi $X$ didefinisikan oleh
$$m(i) = {1\over {2^i}}\ .$$ (Ini tidak lain adalah distribusi geometrik dengan
parameter $1/2$.) Dengan demikian, kita memperoleh
\begin{eqnarray*}  E(X) & = &\sum_{i = 1}^\infty i {1\over{2^i}} \\ & = & \sum_{i =
1}^\infty {1\over{2^i}} + \sum_{i = 2}^\infty {1\over{2^i}} + \cdots \\ & = & 1 +
{1\over 2} + {1\over{2^2}} + \cdots \\ & = & 2\ .
\end{eqnarray*}

\end{example}

\begin{example}(Lanjutan Contoh~\ref{exam 6.05})\label{exam 6.055} Misalkan kita
melempar sebuah koin hingga kepala muncul untuk pertama kalinya; jika banyaknya
lemparan sama dengan $n$, kita dibayar $2^n$ dolar. Berapakah nilai harapan pembayaran
itu?
\par Misalkan $Y$ menyatakan pembayaran tersebut. Maka,
$$P(Y = 2^n) = {1\over{2^n}}\ ,$$ untuk $n \ge 1$. Dengan demikian,
$$E(Y) = \sum_{n = 1}^\infty 2^n {1\over{2^n}}\ ,$$ yang merupakan jumlah divergen.
Jadi, $Y$ tidak mempunyai nilai harapan. Contoh ini disebut \emx {Paradoks
St. Petersburg}.\index{Paradoks St. Petersburg}
Fakta bahwa jumlah di atas tak hingga menunjukkan bahwa seorang pemain seharusnya
bersedia membayar berapa pun jumlah tetap per permainan demi hak memainkan permainan
ini. Pembaca diminta mempertimbangkan berapa banyak yang bersedia ia bayar untuk hak
tersebut. Kecil kemungkinan jawaban pembaca melebihi 10 dolar; di situlah letak
paradoksnya.
\par Pada masa awal sejarah peluang, berbagai matematikawan mengemukakan cara untuk
menyelesaikan paradoks ini. Salah satu gagasan (oleh G.~Cramer)\index{CRAMER, G.}
adalah menganggap bahwa jumlah uang di dunia terhingga. Jadi, ia menganggap terdapat
suatu nilai tetap $n$ sedemikian sehingga jika banyaknya lemparan sama dengan atau
melebihi $n$, pembayarannya adalah $2^n$ dolar. Dalam Latihan~\ref{exer 6.1.21},
pembaca diminta menunjukkan bahwa nilai harapan pembayaran itu kini terhingga.
\par Daniel Bernoulli\index{BERNOULLI, D.} dan Cramer juga mempertimbangkan cara lain
untuk memberikan nilai pada pembayaran. Gagasan mereka adalah bahwa nilai suatu
pembayaran merupakan fungsi tertentu dari pembayaran itu; fungsi semacam itu kini
disebut fungsi utilitas\index{fungsi utilitas}. Contoh fungsi utilitas yang wajar dapat
berupa fungsi akar kuadrat atau fungsi logaritma. Dalam kedua kasus, nilai $2n$ dolar
lebih kecil daripada dua kali nilai $n$ dolar. Dapat ditunjukkan dengan mudah bahwa
dalam kedua kasus, nilai harapan utilitas pembayaran itu terhingga (lihat
Latihan~\ref{exer 6.1.21}).
\end{example}

\begin{example}\label{exam 6.8} Misalkan $T$ adalah waktu sampai keberhasilan
pertama dalam suatu proses percobaan Bernoulli. Kita ambil sebagai ruang sampel
$\Omega$ bilangan bulat $1,~2,~\ldots\ $ dan menetapkan distribusi geometrik
$$ m(j) = P(T = j) = q^{j - 1}p\ . $$ Dengan demikian,
\begin{eqnarray*}  E(T) & = & 1 \cdot p + 2qp + 3q^2p +\cdots \\
     & = & p(1 + 2q + 3q^2 +\cdots )\ .
\end{eqnarray*} Sekarang, jika $|x| < 1$, maka
$$ 1 + x + x^2 + x^3 + \cdots = \frac 1{1 - x}\ .$$ Dengan mendiferensialkan rumus
ini, kita memperoleh
$$ 1 + 2x + 3x^2 +\cdots = \frac 1{(1 - x)^2}\ ,$$ sehingga
$$ E(T) = \frac p{(1 - q)^2} = \frac p{p^2} = \frac 1p\ .$$ Secara khusus, kita
melihat bahwa jika sebuah koin adil dilempar berulang kali, nilai harapan waktu sampai
kepala pertama adalah 1/(1/2) = 2. Jika sebuah dadu dilempar berulang kali, nilai
harapan banyaknya lemparan sampai enam pertama adalah 1/(1/6) = 6.
\end{example}

\subsection*{Penafsiran Nilai Harapan}

Dalam statistika, kita sering berkepentingan dengan nilai rata-rata suatu kumpulan
data. Contoh berikut menunjukkan bahwa gagasan nilai rata-rata dan nilai harapan
berkaitan sangat erat.

\begin{example}\label{exam 6.8.5} Tinggi badan, dalam inci, para pemain perempuan
dalam tim bola basket Swarthmore adalah 5' 9", 5' 9", 5' 6", 5' 8", 5' 11",
5' 5", 5' 7", 5' 6", 5' 6", 5' 7", 5' 10", dan 6' 0".
\par Seorang statistikawan akan menghitung rata-rata tinggi badan (dalam inci) sebagai
berikut:
$$\frac{69 + 69 + 66 + 68 + 71 + 65 + 67 + 66 + 66 + 67 + 70 + 72}{12} = 67.9\ .$$
Bilangan ini juga dapat ditafsirkan sebagai nilai harapan suatu peubah acak. Untuk
melihatnya, misalkan suatu percobaan terdiri atas pemilihan salah seorang pemain itu
secara acak, dan misalkan $X$ menyatakan tinggi badannya. Maka nilai harapan $X$ sama
dengan 67.9.
\end{example}

\par Tentu saja, sebagaimana pada penafsiran frekuensi atas peluang, penafsiran nilai
harapan sebagai hasil rata-rata memerlukan pembenaran lebih lanjut. Kita tahu bahwa
untuk setiap percobaan terhingga, rata-rata hasilnya tidak dapat diramalkan. Namun,
kelak kita akan membuktikan bahwa rata-rata itu biasanya dekat dengan $E(X)$ jika
percobaan diulang sangat banyak kali. Mula-mula kita perlu mengembangkan beberapa
sifat nilai harapan. Dengan menggunakan sifat-sifat ini, beserta sifat konsep varians
yang akan diperkenalkan pada bagian berikutnya, kita akan dapat membuktikan
\emx {Hukum Bilangan Besar.} Teorema ini akan membenarkan secara matematis baik konsep
frekuensi peluang kita maupun penafsiran nilai harapan sebagai nilai rata-rata yang
diharapkan dalam sejumlah besar percobaan.


\subsection*{Nilai Harapan Fungsi Peubah Acak}

Misalkan $X$ adalah peubah acak diskret dengan ruang sampel $\Omega$, dan $\phi(x)$
adalah fungsi bernilai real dengan domain $\Omega$. Maka $\phi(X)$ adalah peubah acak
bernilai real. Salah satu cara menentukan nilai harapan $\phi(X)$ adalah terlebih
dahulu menentukan fungsi distribusi peubah acak ini, lalu menggunakan definisi nilai
harapan. Namun, terdapat cara yang lebih baik untuk menghitung nilai harapan
$\phi(X)$, sebagaimana ditunjukkan dalam contoh berikut.

\begin{example}\label{exam 6.1.5} Misalkan sebuah koin dilempar 9 kali, dengan hasil
\[ HHHTTTTHT\ . \]    
\noindent Kelompok pertama yang terdiri atas tiga kepala disebut suatu
\emx {rentetan}\index{rentetan}. Dalam barisan ini terdapat tiga rentetan lagi, yaitu
empat ekor berikutnya, kepala berikutnya, dan ekor berikutnya. Kita tidak menganggap
dua lemparan pertama saja membentuk suatu rentetan, karena lemparan ketiga mempunyai
nilai yang sama dengan kedua lemparan pertama.
\par Sekarang misalkan suatu percobaan terdiri atas pelemparan sebuah koin adil tiga
kali. Carilah nilai harapan banyaknya rentetan. Akan berguna jika kita memikirkan dua
peubah acak, $X$ dan $Y$, yang berkaitan dengan percobaan ini. Misalkan $X$ menyatakan
barisan kepala dan ekor yang dihasilkan ketika percobaan dilakukan, dan $Y$ menyatakan
banyaknya rentetan dalam hasil $X$. Hasil-hasil $X$ yang mungkin dan nilai-nilai $Y$
yang bersesuaian ditampilkan dalam Tabel~\ref{table 6.2}.
\begin{table}
\centering
$$\begin{tabular}{cc} X    &  Y \\ \hline HHH  & 1\\ HHT  & 2\\ HTH  & 3\\ HTT  & 2\\
THH  & 2\\ THT  & 3\\ TTH  & 2\\ TTT  & 1\\
\end{tabular}
$$
\caption{Melempar sebuah koin tiga kali.}
\label{table 6.2}
\end{table}

Untuk menghitung $E(Y)$ dengan definisi nilai harapan, mula-mula kita harus mencari
fungsi distribusi $m(y)$ dari $Y$; artinya, kita mengelompokkan nilai-nilai $X$ yang
mempunyai nilai $Y$ yang sama, lalu menjumlahkan peluangnya. Dalam kasus ini, kita
menghitung bahwa fungsi distribusi $Y$ adalah: $m(1) = 1/4,\ m(2) = 1/2,$ dan
$m(3) = 1/4$. Dengan mudah diperoleh $E(Y) = 2$.
\par Sekarang misalkan kita tidak mengelompokkan nilai-nilai $X$ yang mempunyai nilai
$Y$ yang sama. Sebaliknya, untuk setiap nilai $X$, yaitu $x$, kita mengalikan peluang
$x$ dengan nilai $Y$ yang bersesuaian, lalu menjumlahkan hasilnya. Kita memperoleh
$$1\biggl(\frac 18\biggr) +2\biggl(\frac 18\biggr) +3\biggl(\frac 18\biggr)
+2\biggl(\frac 18\biggr) +2\biggl(\frac 18\biggr) +3\biggl(\frac 18\biggr)
+2\biggl(\frac 18\biggr) +1\biggl(\frac 18\biggr)\ ,$$ yang sama dengan 2.
\par Hal ini menggambarkan asas umum berikut. Jika $X$ dan $Y$ adalah dua peubah acak,
dan $Y$ dapat ditulis sebagai fungsi $X$, maka nilai harapan $Y$ dapat dihitung dengan
menggunakan fungsi distribusi $X$.
\end{example}

\begin{theorem}\label{thm 6.3.5} Jika $X$ adalah peubah acak diskret dengan ruang
sampel $\Omega$ dan fungsi distribusi $m(x)$, dan jika
%$\phi : \Omega \to {\bf\rm R}$ is a function, then 
$\phi : \Omega \to$ \mat{\rm R} adalah suatu fungsi, maka
$$ E(\phi(X)) = \sum_{x \in \Omega} \phi(x) m(x)\ ,
$$ asalkan deret tersebut konvergen mutlak.
\end{theorem}
 
Bukti teorema ini langsung, tidak lebih dari mengelompokkan nilai-nilai $X$ yang
mempunyai nilai $Y$ yang sama, seperti dalam Contoh~\ref{exam 6.1.5}.

\subsection*{Jumlah Dua Peubah Acak}

Banyak hasil penting dalam teori peluang berkaitan dengan jumlah peubah acak. Mula-mula
kita tinjau apa artinya menjumlahkan dua peubah acak.

\begin{example}\label{exam 6.06} Kita melempar sebuah koin dan menetapkan $X$ bernilai
1 jika koin menghasilkan kepala dan 0 jika koin menghasilkan ekor. Kemudian kita
melempar sebuah dadu dan misalkan $Y$ menyatakan sisi yang muncul. Apa arti $X+Y$,
dan bagaimana distribusinya? Dalam kasus ini pertanyaan tersebut mudah dijawab dengan
meninjau, seperti yang kita lakukan dalam Bab~\ref{chp 4}, peubah acak gabungan
$Z = (X,Y)$, yang hasil-hasilnya berupa pasangan terurut $(x, y)$, dengan
$0 \le x \le 1$ dan $1 \le y \le 6$. Deskripsi percobaan itu membuat kita wajar
menganggap bahwa $X$ dan $Y$ saling bebas, sehingga fungsi distribusi $Z$ seragam,
dengan $1/12$ ditetapkan bagi setiap hasil. Sekarang himpunan hasil $X+Y$ beserta
fungsi distribusinya mudah ditemukan.
\end{example}

Dalam Contoh~\ref{exam 6.03}, peubah acak $X$ menyatakan banyaknya kepala yang muncul
ketika sebuah koin adil dilempar tiga kali. Wajar untuk memandang $X$ sebagai jumlah
peubah-peubah acak $X_1, X_2, X_3$, dengan $X_i$ didefinisikan bernilai 1 jika
lemparan ke-$i$ menghasilkan kepala, dan 0 jika lemparan ke-$i$ menghasilkan ekor.
Nilai harapan setiap $X_i$ sangat mudah dihitung. Ternyata nilai harapan $X$ dapat
diperoleh hanya dengan menjumlahkan nilai harapan setiap $X_i$. Fakta ini dinyatakan
dalam teorema berikut.

\begin{theorem}\label{thm 6.1} Misalkan $X$ dan $Y$ adalah peubah acak dengan nilai
harapan terhingga. Maka
$$ E(X + Y) = E(X) + E(Y)\ ,
$$ dan jika $c$ adalah sebarang konstanta, maka
$$ E(cX) = cE(X)\ .
$$
\proof Misalkan ruang sampel $X$ dan $Y$ masing-masing dilambangkan dengan $\Omega_X$
dan $\Omega_Y$, serta misalkan
$$\Omega_X = \{x_1, x_2, \ldots\}$$ dan
$$\Omega_Y = \{y_1, y_2, \ldots\}\ .$$ Maka peubah acak $X + Y$ dapat kita pandang
sebagai hasil penerapan fungsi $\phi(x, y) = x + y$ pada peubah acak gabungan
$(X,Y)$. Dengan Teorema~\ref{thm 6.3.5}, kita memperoleh
\begin{eqnarray*} E(X+Y) & = &\sum_j \sum_k (x_j + y_k) P(X = x_j,\ Y = y_k) \\
       & = &\sum_j \sum_k x_j P(X = x_j,\ Y = y_k) +
\sum_j \sum_k y_k P(X = x_j,\ Y = y_k) \\ & = &\sum_j x_j P(X = x_j) + \sum_k y_k P(Y
= y_k)\ .
\end{eqnarray*} Kesamaan terakhir mengikuti fakta bahwa
$$\sum_k P(X = x_j,\ Y = y_k)\ \  =\ \  P(X = x_j)$$ dan
$$\sum_j P(X = x_j,\ Y = y_k)\ \  =\ \  P(Y = y_k)\ .$$ Dengan demikian,
$$E(X+Y) = E(X) + E(Y)\ .$$ Jika $c$ adalah sebarang konstanta,
\begin{eqnarray*}  E(cX) & = & \sum_j cx_j P(X = x_j) \\
      & = & c\sum_j x_j P(X = x_j)\\
      & = & cE(X)\ .
\end{eqnarray*}
\end{theorem}
Dapat dibuktikan dengan mudah melalui induksi matematika bahwa \emx {nilai harapan
jumlah sebarang banyak terhingga peubah acak sama dengan jumlah nilai harapan setiap
peubah acak tersebut.}
\par Penting diperhatikan bahwa kebebasan bersama suku-suku penjumlah tidak diperlukan
sebagai hipotesis dalam Teorema~\ref{thm 6.1} maupun generalisasinya. Fakta bahwa nilai
harapan dapat dijumlahkan, terlepas dari apakah suku-suku penjumlah saling bebas,
kadang-kadang disebut Misteri Fundamental Pertama Peluang.\index{Misteri Fundamental
Pertama Peluang}

\begin{example}\label{exam 6.1} Misalkan $Y$ adalah banyaknya titik tetap dalam suatu
permutasi acak himpunan $\{a,b,c\}$. Untuk mencari nilai harapan $Y$, akan berguna
jika kita meninjau peubah acak dasar yang berkaitan dengan percobaan ini, yaitu peubah
acak $X$ yang menyatakan permutasi acak tersebut. Ada enam hasil $X$ yang mungkin,
dan masing-masing kita beri peluang $1/6$ (lihat Tabel~\ref{table 6.3}). Kemudian kita
dapat menghitung $E(Y)$ dengan Teorema~\ref{thm 6.3.5}, sebagai
$$3\Bigl({1\over 6}\Bigr) + 1\Bigl({1\over 6}\Bigr) + 1\Bigl({1\over 6}\Bigr) +
0\Bigl({1\over 6}\Bigr) +  0\Bigl({1\over 6}\Bigr) + 1\Bigl({1\over 6}\Bigr) = 1\ .
$$
\begin{table}
\centering
\begin{tabular}{ccc} &                          & \\ 
\hline &$X$                 & $Y$ \\
\hline  &$a\;\;\;b\;\;\;   c$       & 3 \\  &$a\;\;\; c\;\;\;  b$       & 1 \\ 
&$b\;\;\; a\;\;\;  c$       & 1 \\  &$b\;\;\; c\;\;\;  a$       & 0 \\  &$c\;\;\;
a\;\;\;  b$       & 0 \\  &$c\;\;\; b\;\;\;  a$       & 1 \\ 
\hline
\end{tabular}
\caption{Banyaknya titik tetap.}
\label{table 6.3}
\end{table}

\par Sekarang kita berikan cara yang sangat cepat untuk menghitung rata-rata banyaknya
titik tetap dalam suatu permutasi acak himpunan $\{1, 2, 3, \ldots, n\}$. Misalkan
$Z$ menyatakan permutasi acak tersebut. Untuk setiap $i$, $1 \le i \le n$, misalkan
$X_i$ sama dengan 1 jika $Z$ menetapkan $i$, dan 0 jika tidak. Jadi, jika $F$
menyatakan banyaknya titik tetap dalam $Z$, maka
$$F = X_1 + X_2 + \cdots + X_n\ .$$ Karena itu, Teorema~\ref{thm 6.1} menyiratkan
$$E(F) = E(X_1) + E(X_2) + \cdots + E(X_n)\ .$$ Namun, mudah dilihat bahwa untuk
setiap $i$,
$$E(X_i) = {1\over n}\ ,$$ sehingga
$$E(F) = 1\ .$$ Metode penghitungan nilai harapan ini sering sangat berguna. Metode
tersebut berlaku setiap kali peubah acak yang ditinjau dapat ditulis sebagai jumlah
peubah-peubah acak yang lebih sederhana. Sekali lagi kita tekankan bahwa suku-suku
penjumlah itu tidak harus saling bebas.
\end{example}

\subsection*{Percobaan Bernoulli}

\begin{theorem}\label{thm 6.3} Misalkan $S_n$ adalah banyaknya keberhasilan dalam
$n$ percobaan Bernoulli dengan peluang keberhasilan $p$ pada setiap percobaan. Maka
nilai harapan banyaknya keberhasilan adalah $np$. Artinya,
$$ E(S_n) = np\ .
$$
\proof Misalkan $X_j$ adalah peubah acak yang bernilai~1 jika hasil ke-$j$ merupakan
keberhasilan dan~0 jika merupakan kegagalan. Maka, untuk setiap $X_j$,
$$ E(X_j) = 0\cdot(1 - p) + 1\cdot p = p\ .
$$ Karena
$$ S_n = X_1 + X_2 +\cdots+ X_n\ ,
$$ dan nilai harapan suatu jumlah sama dengan jumlah nilai-nilai harapannya, kita
memperoleh
\begin{eqnarray*} E(S_n) & = & E(X_1) + E(X_2) +\cdots+ E(X_n) \\
       & = & np\ .
\end{eqnarray*}
\end{theorem}

\subsection*{Distribusi Poisson}

Ingat bahwa distribusi Poisson dengan parameter $\lambda$ diperoleh sebagai limit
distribusi binomial dengan parameter $n$ dan $p$, dengan asumsi $np = \lambda$ dan
$n \rightarrow \infty$. Karena untuk setiap $n$ distribusi binomial yang bersesuaian
mempunyai nilai harapan $\lambda$, wajar untuk menduga bahwa distribusi Poisson dengan
parameter $\lambda$ juga mempunyai nilai harapan $\lambda$. Memang demikian halnya,
dan pembaca dipersilakan menunjukkannya (lihat Latihan~\ref{exer 6.1.100}).

\subsection*{Kebebasan}

Jika $X$ dan $Y$ adalah dua peubah acak, secara umum tidak benar bahwa $E(X
\cdot Y) = E(X)E(Y)$. Namun, hal ini benar jika $X$ dan $Y$ \emx {saling bebas.}

\begin{theorem}\label{thm 6.4} Jika $X$ dan $Y$ adalah peubah acak yang saling bebas,
maka
$$ E(X \cdot Y) = E(X)E(Y)\ .
$$
\proof Misalkan
$$\Omega_X = \{x_1, x_2, \ldots\}$$ dan
$$\Omega_Y = \{y_1, y_2, \ldots\}$$ masing-masing adalah ruang sampel $X$ dan $Y$.
Dengan menggunakan Teorema~\ref{thm 6.3.5}, kita memperoleh
$$ E(X \cdot Y) = \sum_j \sum_k x_jy_k P(X = x_j,\ Y = y_k)\ .
$$

Namun, jika $X$ dan $Y$ saling bebas,
$$ P(X = x_j, Y = y_k) = P(X = x_j)P(Y = y_k)\ .
$$ Dengan demikian,
\begin{eqnarray*} E(X \cdot Y) & = & \sum_j\sum_k x_j y_k P(X = x_j) P(Y = y_k) \\
             & = & \left(\sum_j x_j P(X = x_j)\right) \left(\sum_k y_k P(Y =
y_k)\right) \\
             & = &E(X) E(Y)\ .
\end{eqnarray*}
\end{theorem}

\begin{example}\label{exam 6.3} Sebuah koin dilempar dua kali. $X_i = 1$ jika
lemparan ke-$i$ menghasilkan kepala dan~0 jika tidak. Kita tahu bahwa $X_1$ dan $X_2$
saling bebas. Masing-masing mempunyai nilai harapan 1/2. Jadi,
$E(X_1 \cdot X_2) = E(X_1) E(X_2) = (1/2)(1/2) = 1/4$.
\end{example}

Berikutnya kita berikan contoh sederhana untuk menunjukkan bahwa nilai harapan tidak
harus dapat dikalikan jika peubah-peubah acaknya tidak saling bebas.

\begin{example}\label{exam 6.4} Tinjaulah satu kali pelemparan koin. Kita definisikan
peubah acak $X$ bernilai~1 jika muncul kepala dan~0 jika muncul ekor, serta kita
tetapkan $Y = 1 - X$. Maka $E(X) = E(Y) = 1/2$. Namun, $X \cdot Y = 0$ untuk kedua
hasil. Karena itu, $E(X \cdot Y) = 0 \ne E(X) E(Y)$.
\end{example}

Kita kembali ke contoh rekor dalam Bagian~\ref{sec 3.1} untuk penerapan lain dari
hasil bahwa nilai harapan jumlah peubah acak sama dengan jumlah nilai harapan setiap
peubah acak tersebut.

\subsection*{Rekor}\index{rekor}

\begin{example}\label{exam 6.5} Kita mulai mencatat rekor curah salju tahun ini dan
ingin mencari nilai harapan banyaknya rekor yang akan terjadi dalam $n$ tahun
berikutnya. Tahun pertama pasti merupakan rekor. Tahun kedua akan menjadi rekor jika
curah salju pada tahun kedua lebih besar daripada pada tahun pertama. Berdasarkan
simetri, peluang ini adalah 1/2. Secara lebih umum, misalkan $X_j$ bernilai~1 jika
tahun ke-$j$ merupakan rekor dan~0 jika tidak. Untuk mencari $E(X_j)$, kita hanya
perlu mencari peluang bahwa tahun ke-$j$ merupakan rekor. Namun, rekor curah salju
selama $j$ tahun pertama mempunyai kemungkinan yang sama untuk jatuh pada salah satu
dari tahun-tahun tersebut, sehingga $E(X_j) = 1/j$. Oleh karena itu, jika $S_n$
adalah jumlah seluruh rekor yang diamati selama $n$ tahun pertama,
$$ E(S_n) = 1 + \frac 12 + \frac 13 +\cdots+ \frac 1n\ .
$$ Ini adalah \emx {deret harmonik divergen} yang terkenal. Mudah ditunjukkan bahwa
$$ E(S_n) \sim \log n
$$ ketika $n \rightarrow \infty$. Hampiran yang lebih akurat bagi $E(S_n)$ diberikan
oleh bentuk
$$\log n + \gamma + {1\over {2n}}\ ,$$
dengan $\gamma$ menyatakan konstanta Euler, yang kira-kira sama dengan .5772.

Jadi, dalam sepuluh tahun nilai harapan banyaknya rekor kira-kira $2.9298$; nilai
eksaknya adalah jumlah sepuluh suku pertama deret harmonik, yaitu 2.9290.
\end{example}

\subsection*{Craps}\index{craps}

\begin{example}\label{exam 6.6} Dalam permainan craps, pemain memasang taruhan dan
melempar sepasang dadu. Jika jumlah kedua angka adalah 7~atau~11, pemain menang; jika
jumlahnya 2,~3, atau~12, pemain kalah. Jika diperoleh bilangan lain, misalnya $r$,
maka $r$ menjadi poin pemain dan ia terus melempar hingga muncul $r$ atau 7. Jika $r$
muncul lebih dahulu, ia menang; jika 7 muncul lebih dahulu, ia kalah. Program
{\bf Craps}\index{Craps (program)} menyimulasikan permainan ini beberapa kali.
\par Kita telah menjalankan program untuk 1000 permainan, dengan pemain memasang
taruhan 1~dolar setiap kali. Rata-rata kemenangan pemain adalah $-.006$. Tampaknya
permainan craps hanya sedikit merugikan. Mari kita hitung nilai harapan kemenangan
dalam satu permainan untuk memeriksa apakah memang demikian. Kita menyusun ukuran
pohon dua tahap seperti ditunjukkan dalam Gambar~\ref{fig 6.1}.

\putfig{4.5truein}{PSfig6-1}{Ukuran pohon untuk craps.}{fig 6.1}

Tahap pertama menyatakan jumlah-jumlah yang mungkin pada lemparan pertamanya. Tahap
kedua menyatakan hasil-hasil permainan yang mungkin jika permainan belum berakhir
pada lemparan pertama. Pada tahap ini kita menyatakan hasil-hasil yang mungkin dari
suatu barisan lemparan yang diperlukan untuk menentukan hasil akhir. Peluang cabang
untuk tahap pertama dihitung dengan cara biasa, dengan menganggap semua 36 hasil yang
mungkin bagi pasangan dadu mempunyai peluang yang sama. Untuk tahap kedua, kita
menganggap permainan pada akhirnya akan berakhir, lalu menghitung peluang bersyarat
untuk memperoleh poin atau~7. Sebagai contoh, anggaplah poin pemain adalah~6. Maka
permainan akan berakhir ketika salah satu dari sebelas pasangan, $(1,5)$, $(2,4)$,
$(3,3)$, $(4,2)$, $(5,1)$,
$(1,6)$, $(2,5)$,
$(3,4)$, $(4,3)$, $(5,2)$, $(6,1)$, muncul. Kita menganggap setiap pasangan yang
mungkin ini mempunyai peluang yang sama. Pemain menang dalam lima kasus pertama dan
kalah dalam enam kasus terakhir. Jadi, peluang menang adalah 5/11 dan peluang kalah
adalah 6/11. Dari peluang lintasan, kita dapat mencari peluang bahwa pemain menang
1~dolar, yaitu 244/495. Peluang kalah kemudian 251/495. Jadi, jika $X$ adalah
kemenangannya untuk taruhan satu dolar,
\begin{eqnarray*} E(X) & = & 1\Bigl(\frac {244}{495}\Bigr) + (-1)\Bigl(\frac
{251}{495}\Bigr) \\
     & = & -\frac {7}{495} \approx -.0141\ .
\end{eqnarray*} Permainan tersebut merugikan, tetapi hanya sedikit. Nilai harapan
keuntungan pemain dalam $n$ permainan adalah $-n(.0141)$. Jika $n$ tidak besar, ini
merupakan nilai harapan kerugian yang kecil bagi pemain. Kasino menjalankan sangat
banyak permainan, sehingga mampu memperoleh rata-rata keuntungan kecil per permainan
namun tetap mengharapkan laba besar.
\end{example}

\subsection*{Roulette}\index{roulette}

\begin{example}\label{exam 6.7} Di Las Vegas, sebuah roda roulette mempunyai 38 petak
bernomor 0,\ 00,\ 1,\ 2,\ \ldots,\ 36. Petak 0 dan 00 berwarna hijau, sedangkan
separuh dari 36 petak sisanya berwarna merah dan separuh lainnya hitam. Seorang bandar
memutar roda dan melemparkan bola gading. Jika Anda bertaruh 1 dolar pada merah, Anda
menang 1 dolar jika bola berhenti pada petak merah; jika tidak, Anda kehilangan satu
dolar. Kita ingin menghitung nilai harapan kemenangan Anda jika Anda bertaruh 1 dolar
pada merah.
\par
Misalkan $X$ adalah peubah acak yang menyatakan kemenangan Anda dalam taruhan 1 dolar
pada merah dalam roulette Las Vegas. Distribusi $X$ diberikan oleh
$$ m_{X} = \pmatrix{
   -1 & 1 \cr 20/38 & 18/38 \cr},
$$ dan dapat dihitung dengan mudah (lihat Latihan~\ref{exer 6.1.5}) bahwa
$$ 
E(X) \approx -.0526\ .
$$
\par
Sekarang kita tinjau permainan roulette di Monte Carlo, dengan mengikuti pembahasan
Sagan.\footnote{H. Sagan, \emx {Markov Chains in Monte Carlo,} Math. Mag., vol. 54, no. 1 (1981),
pp. 3-10.}\index{SAGAN, H.} Dalam permainan roulette di Monte Carlo hanya terdapat
satu~0. Jika Anda bertaruh 1 franc pada merah dan muncul~0, maka, bergantung pada
kasinonya, satu atau lebih pilihan berikut dapat ditawarkan:
\par
\noindent
(a) Anda menerima kembali 1/2 taruhan Anda, dan kasino mengambil separuh lainnya.
\par
\noindent
(b) Taruhan Anda ``dipenjara," yang akan kita lambangkan dengan $P_1$. Jika merah
muncul pada putaran berikutnya, Anda menerima kembali taruhan Anda (tetapi tidak
memenangkan uang). Jika hitam atau 0 muncul, Anda kehilangan taruhan Anda.
\par
\noindent
(c) Taruhan Anda dipenjara dalam $P_1$, seperti sebelumnya. Jika merah muncul pada
putaran berikutnya, taruhan Anda dikembalikan; jika hitam muncul pada putaran
berikutnya, Anda kehilangan taruhan. Jika 0 muncul pada putaran berikutnya, taruhan
Anda dimasukkan ke penjara ganda, yang akan kita lambangkan dengan $P_2$. Jika taruhan
Anda berada dalam penjara ganda dan merah muncul pada putaran berikutnya, taruhan itu
dipindahkan kembali ke penjara $P_1$ dan permainan berlanjut seperti sebelumnya. Jika
taruhan Anda berada dalam penjara ganda dan hitam atau 0 muncul pada putaran
berikutnya, Anda kehilangan taruhan. Lihat Gambar~\ref{fig 6.1.5}, yang menampilkan
pohon untuk pilihan ini. Dalam gambar tersebut, $S$ adalah posisi awal, $W$ berarti
Anda memenangkan taruhan, $L$ berarti Anda kehilangan taruhan, dan $E$ berarti Anda
impas.
\putfig{4.5truein}{PSfig6-1-5}{Pohon untuk roulette Monte Carlo dengan 2 penjara.}{fig 6.1.5}
\par
Menarik untuk membandingkan nilai harapan kemenangan dari taruhan 1 franc pada merah
di bawah masing-masing dari ketiga pilihan ini. Dua perhitungan pertama kita serahkan
sebagai latihan (lihat Latihan~\ref{exer 6.1.38}). Misalkan Anda memilih alternatif
(c). Perhitungan untuk kasus ini menggambarkan cara para ahli peluang Prancis pada masa
awal mengerjakan masalah seperti ini.
\par
Misalkan Anda bertaruh pada merah, memilih alternatif (c), dan muncul~0. Hasil-hasil
mendatang yang mungkin ditampilkan dalam diagram pohon pada Gambar~\ref{fig 6.2}.
Anggaplah uang Anda berada di penjara pertama, dan misalkan $x$ adalah peluang Anda
kehilangan franc tersebut. Dari diagram pohon kita melihat bahwa
$$
x = \frac {18}{37} + \frac 1{37}P({\rm Anda\ kehilangan\ franc\ Anda\ }|\ {\rm franc\ Anda\ berada
\ di\ }P_2)\ .
$$
Selain itu,
$$
P({\rm Anda\ kehilangan\ franc\ Anda\ }|\ {\rm franc\ Anda\ berada
\ di\ }P_2) =  \frac {19}{37} + \frac{18}{37}x\ .
$$
Jadi, kita mempunyai
$$x = \frac{18}{37} + \frac 1{37}\Bigl(\frac {19}{37} + \frac{18}{37}x\Bigr)\ .$$
Dengan menyelesaikan persamaan bagi $x$, kita memperoleh $x = 685/1351$. Jadi, jika
dimulai dari $S$, peluang bahwa Anda kehilangan taruhan sama dengan
$$\frac {18}{37} + \frac 1{37}x = \frac{25003}{49987}\ .$$
\par
Untuk mencari peluang bahwa Anda menang ketika bertaruh pada merah, perhatikan bahwa
Anda hanya dapat menang jika merah muncul pada putaran pertama, dan hal ini terjadi
dengan peluang 18/37. Dengan demikian, nilai harapan kemenangan Anda adalah
$$ 1 \cdot {\frac{18}{37}} -1 \cdot {\frac {25003}{49987}} = -\frac{687}{49987} 
\approx -.0137\ .$$

\putfig{4.5truein}{PSfig6-2}{Uang Anda dimasukkan ke penjara.}  {fig 6.2}

Menarik untuk diperhatikan bahwa pilihan (c) yang lebih romantis ternyata kurang
menguntungkan daripada pilihan (a) (lihat Latihan~\ref{exer 6.1.38}).
\par
Jika Anda bertaruh 1~dolar pada bilangan 17, fungsi distribusi kemenangan Anda $X$
adalah
$$ P_X = \pmatrix{
   -1 & 35 \cr 36/37 & 1/37 \cr}\ ,
$$ 
dan nilai harapan kemenangannya adalah
$$ -1 \cdot {\frac{36}{37}} + 35 \cdot {\frac 1{37}} = -\frac 1{37} \approx -.027\ .
$$ 
Jadi, di Monte Carlo taruhan yang berbeda mempunyai nilai harapan yang berbeda. Di
Las Vegas, hampir semua taruhan mempunyai nilai harapan yang sama, yaitu
$-2/38 = -.0526$ (lihat Latihan~\ref{exer 6.1.4} dan \ref{exer 6.1.5}).
\end{example}

\subsection*{Nilai Harapan Bersyarat}

\begin{definition}\label{def 6.3} Jika $F$ adalah sebarang kejadian dan $X$ adalah
peubah acak dengan ruang sampel $\Omega = \{x_1, x_2, \ldots\}$, maka
\emx {nilai harapan bersyarat\index{nilai harapan bersyarat} dengan syarat $F$}
didefinisikan oleh
$$ E(X|F) = \sum_j x_j P(X = x_j|F)\ .
$$ Nilai harapan bersyarat paling sering digunakan dalam bentuk yang diberikan oleh
teorema berikut.
\end{definition}

\begin{theorem}\label{thm 6.5} Misalkan $X$ adalah peubah acak dengan ruang sampel
$\Omega$. Jika $F_1$,~$F_2$, \dots,~$F_r$ adalah kejadian-kejadian sedemikian sehingga
$F_i \cap F_j = \emptyset$ untuk $i \ne j$ dan $\Omega = \cup_j F_j$, maka
$$ E(X) = \sum_j E(X|F_j) P(F_j)\ .
$$
\proof Kita mempunyai
\begin{eqnarray*}
\sum_j E(X|F_j) P(F_j) & = & \sum_j \sum_k x_k P(X = x_k|F_j) P(F_j) \\
                       & = & \sum_j \sum_k x_k P(X = x_k\,\, {\rm dan}\,\, 
F_j\,\,{\rm terjadi}) \\
                       & = & \sum_k \sum_j x_k P(X = x_k\,\,{\rm dan}\,\,F_j\,\,{\rm
terjadi})
\\
                       & = & \sum_k x_k P(X = x_k) \\
                       & = & E(X)\ .
\end{eqnarray*}
\end{theorem}

\begin{example}(Lanjutan Contoh~\ref{exam 6.6}){\label{exam 6.10}}\index{craps}
Misalkan $T$ adalah banyaknya lemparan dalam satu permainan craps. Kita dapat memandang
satu permainan sebagai proses dua tahap. Tahap pertama terdiri atas satu lemparan
sepasang dadu. Permainan berakhir jika lemparan ini menghasilkan 2, 3, 7, 11, atau 12.
Jika tidak, poin pemain ditetapkan dan tahap kedua dimulai. Tahap kedua ini terdiri atas
barisan lemparan yang berakhir ketika poin pemain atau 7 muncul. Kita mencatat hasil
percobaan dua tahap ini dengan menggunakan peubah acak $X$ dan $S$, dengan $X$
menyatakan lemparan pertama dan $S$ menyatakan banyaknya lemparan pada tahap kedua
percobaan (tentu saja, kadang-kadang $S$ sama dengan 0). Perhatikan bahwa $T = S+1$.
Kemudian, berdasarkan Teorema~\ref{thm 6.5},
$$ E(T) = \sum_{j = 2}^{12} E(T|X = j) P(X = j)\ .
$$ Jika $j = 7$,~11 atau~2, 3,~12, maka $E(T|X = j) = 1$. Jika $j = 4, 5, 6, 8,
9,$ atau $10$, kita dapat menggunakan Contoh~\ref{exam 6.8} untuk menghitung nilai
harapan $S$. Dalam setiap kasus ini, kita terus melempar hingga memperoleh $j$ atau 7.
Jadi, $S$ berdistribusi geometrik dengan parameter $p$, yang bergantung pada $j$.
Sebagai contoh, jika $j = 4$, nilai $p$ adalah $3/36 + 6/36 = 1/4$. Maka, dalam
kasus ini, nilai harapan banyaknya lemparan tambahan adalah $1/p = 4$, sehingga
$E(T|X = 4) = 1 + 4 = 5$. Dengan melakukan perhitungan yang bersesuaian bagi nilai
$j$ lain yang mungkin dan menggunakan Teorema~\ref{thm 6.5}, diperoleh
\begin{eqnarray*} E(T) & = & 1\Bigl(\frac {12}{36}\Bigr)   
    + \Bigl(1 + \frac {36}{3 + 6}\Bigr)\Bigl(\frac 3{36}\Bigr)   
    + \Bigl(1 + \frac {36}{4 + 6}\Bigr)\Bigl(\frac 4{36}\Bigr) \\
& & + \Bigl(1 + \frac {36}{5 + 6}\Bigr)\Bigl(\frac 5{36}\Bigr)  
    + \Bigl(1 + \frac {36}{5 + 6}\Bigr)\Bigl(\frac 5{36}\Bigr) \\
& & + \Bigl(1 + \frac {36}{4 + 6}\Bigr)\Bigl(\frac 4{36}\Bigr) 
    + \Bigl(1 + \frac {36}{3 + 6}\Bigr)\Bigl(\frac 3{36}\Bigr) \\ 
& = & \frac {557}{165} \\
& \approx & 3.375\dots\ .
\end{eqnarray*}
\end{example}

\subsection*{Martingal}

Kita dapat memperluas gagasan keadilan bagi seorang pemain yang memainkan serangkaian
permainan dengan menggunakan konsep nilai harapan bersyarat.

\begin{example}\label{exam 6.11} Misalkan $S_1$,~$S_2$, \dots,~$S_n$ adalah kekayaan
Peter yang terkumpul selama memainkan kepala atau ekor (lihat
Contoh~\ref{exam 1.3}). Maka
$$ E(S_n | S_{n - 1} = a,\dots,S_1 = r) = \frac 12 (a + 1) + \frac 12 (a - 1) = a\ .
$$

Kita perhatikan bahwa nilai harapan kekayaan Peter setelah permainan berikutnya sama
dengan kekayaannya saat ini. Jika hal ini terjadi, kita katakan bahwa permainannya
\emx {adil.}\index{permainan adil} Permainan adil juga disebut
\emx {martingal.}\index{martingal} Jika koin tidak seimbang dan menghasilkan kepala
dengan peluang $p$ serta ekor dengan peluang $q = 1 - p$, maka
$$ E(S_n | S_{n - 1} = a,\dots,S_1 = r) = p (a + 1) + q (a - 1) = a + p - q\ .
$$ Jadi, jika $p < q$, permainan ini merugikan, dan jika $p > q$, permainan ini
menguntungkan.
\end{example}

Jika Anda berada di kasino, Anda akan melihat para pemain menerapkan \emx {sistem}
permainan yang rumit\index{sistem perjudian} untuk mencoba mengubah permainan yang
merugikan menjadi menguntungkan. Dua sistem semacam itu, sistem penggandaan martingal
dan sistem Labouchere yang lebih konservatif, dijelaskan dalam
Latihan~\ref{sec 1.1}.\ref{exer 1.1.9}~dan~\ref{sec 1.1}.\ref{exer 1.1.10}.
Sayangnya, sistem semacam itu bahkan tidak dapat mengubah permainan adil menjadi
permainan yang menguntungkan.

Meskipun demikian, mengembangkan sistem permainan untuk perjudian dan permainan lain,
seperti pasar saham, merupakan kegemaran banyak orang. Kita menutup bagian ini dengan
ilustrasi sederhana mengenai sistem semacam itu.

\subsection*{Harga Saham}\index{harga saham}

\begin{example}\label{exam 6.12} Anggaplah nilai suatu saham naik atau turun 1~dolar
setiap hari, masing-masing dengan peluang 1/2. Model yang disederhanakan ini dapat kita
samakan dengan permainan kepala atau ekor yang sudah kita kenal. Kita menganggap
seorang pembeli, Tn.\ Ace\index{Ace, Tn.}, menerapkan strategi berikut. Ia membeli
saham pada hari pertama dengan harga $V$. Kemudian ia menunggu sampai harga saham naik
satu menjadi $V + 1$, lalu menjualnya. Setelah itu ia terus mengamati saham hingga
harganya turun kembali ke $V$. Ia membeli lagi, menunggu sampai harganya naik menjadi
$V + 1$, lalu menjualnya. Jadi, ia memegang saham selama selang-selang ketika harganya
naik 1~dolar. Dalam setiap selang semacam itu, ia memperoleh laba 1~dolar. Namun, kita
menganggap ia hanya dapat melakukannya selama sejumlah terhingga hari perdagangan.
Dengan demikian, ia dapat merugi jika, dalam selang terakhir ketika ia memegang saham,
harganya tidak kembali naik ke $V + 1$; inilah satu-satunya cara ia dapat merugi.
Dalam Gambar~\ref{fig 6.3}, kita menggambarkan riwayat khas jika Tn.\ Ace harus
berhenti setelah dua puluh hari.
\putfig{4truein}{PSfig6-3}{Sistem Tn.\ Ace.}{fig 6.3}
Tn.\ Ace memegang saham berdasarkan sistemnya selama hari-hari yang ditunjukkan dengan
garis putus-putus. Kita perhatikan bahwa untuk riwayat dalam Gambar~\ref{fig 6.3},
sistemnya memberinya keuntungan bersih 4~dolar.
\par
Kita telah menulis program {\bf StockSystem}\index{StockSystem (program)} untuk
menyimulasikan kekayaan Tn.\ Ace jika ia menggunakan sistemnya selama periode $n$
hari. Jika program ini dijalankan sangat banyak kali, misalnya untuk $n = 20$, kita
mendapati bahwa nilai harapan kemenangannya sangat dekat dengan 0, tetapi peluang ia
berada di depan setelah 20 hari jauh lebih besar daripada 1/2. Untuk nilai $n$ yang
kecil, distribusi tepat kemenangannya dapat dihitung. Distribusi untuk kasus $n = 20$
ditampilkan dalam Gambar~\ref{fig 6.3.1}. Dengan menggunakan distribusi ini, mudah
dihitung bahwa nilai harapan kemenangannya tepat 0. Ini merupakan contoh lain dari
fakta bahwa permainan adil (suatu martingal\index{martingal}) tetap adil di bawah
sistem permainan yang cukup umum.
\putfig{4truein}{PSfig6-3-1}{Distribusi kemenangan untuk $n = 20$.}{fig 6.3.1}
\par
Walaupun nilai harapan kemenangannya adalah 0, peluang Tn.\ Ace berada di depan setelah
20 hari sekitar .610. Jadi, jika model acak sederhana kita benar, ia dapat mengatakan
kepada teman-temannya bahwa sistemnya memberinya peluang lebih besar untuk berada di
depan daripada seseorang yang sekadar membeli saham lalu menyimpannya. Sejumlah
penelitian telah dilakukan untuk menentukan seberapa acak pasar saham itu.
\end{example}


\subsection*{Catatan Sejarah}

%**********Put an account of the St. Petersburg paradox in this section.
Dengan Hukum Bilangan Besar yang memperkuat penafsiran frekuensi atas peluang, wajar
bagi kita untuk membenarkan definisi nilai harapan berdasarkan hasil rata-rata dari
sejumlah besar pengulangan percobaan. Konsep nilai harapan telah digunakan sebelum
didefinisikan secara formal; ketika digunakan, konsep itu tidak dipandang sebagai
nilai rata-rata, melainkan sebagai nilai yang sesuai bagi suatu pertaruhan. Sebagai
contoh, ingatlah dari bagian Catatan Sejarah dalam Bab~\ref{chp 1},
Bagian~\ref{sec 1.2}, cara Pascal mencari nilai suatu seri tiga permainan yang harus
dihentikan sebelum selesai.
\par
Pascal\index{PASCAL, B.} mula-mula mengamati bahwa jika masing-masing pemain hanya
perlu memenangkan satu permainan lagi, taruhan sebesar 64~pistole harus dibagi rata.
Kemudian ia meninjau kasus ketika seorang pemain telah memenangkan dua permainan dan
pemain lainnya satu.

Jika pemain pertama menang pada permainan berikutnya, ia memperoleh 64~pistole; jika
kalah, ia masih memperoleh 32. Karena 32~pistole pertama sudah pasti menjadi haknya dan
sisa 32~pistole bergantung pada pertaruhan seimbang, bagian yang adil baginya ialah
$32+16=48$~pistole, sedangkan pemain lain memperoleh 16.\footnote{Ringkasan berdasarkan
F.~N. David, \emx {Games, Gods and Gambling} (London: Griffin, 1962), p.~231;
ungkapan terjemahan modern tidak diterjemahkan atau direproduksi.}

Perhatikan bahwa Pascal mereduksi masalah itu menjadi pertaruhan simetris, ketika
masing-masing pemain memperoleh jumlah yang sama, dan menganggap jelas bahwa dalam
kasus ini taruhan harus dibagi rata.

Kajian sistematis pertama mengenai nilai harapan muncul dalam buku
Huygens.\index{HUYGENS, C.|(} Seperti Pascal, Huygens mencari nilai suatu pertaruhan
dengan menganggap bahwa jawabannya jelas untuk situasi simetris tertentu, lalu
menggunakannya untuk menurunkan nilai harapan bagi situasi umum. Ia melakukannya
secara bertahap. Proposisi pertamanya adalah

\begin{quote} Prop.~I. Jika saya mengharapkan $a$ atau $b$, yang masing-masing dapat
jatuh kepada saya dengan peluang yang sama, maka nilai Harapan saya adalah
$(a + b)/2$, yaitu setengah dari Jumlah $a$ dan $b$.\footnote{C. Huygens,
\emx {Calculating in Games of Chance,} terjemahan yang dinisbatkan kepada John
Arbuthnot (London, 1692), p.~34.}
\end{quote}

Huygens membuktikannya sebagai berikut. Anggaplah dua pemain A dan B memainkan suatu
permainan; masing-masing memasang taruhan $(a + b)/2$ dengan peluang yang sama untuk
memenangkan seluruh taruhan. Maka nilai permainan bagi masing-masing pemain adalah
$(a + b)/2$. Sebagai contoh, jika permainan harus dihentikan, jelas masing-masing
pemain cukup menerima kembali taruhan awalnya. Sekarang, berdasarkan simetri, nilai
ini tidak berubah jika kita menambahkan syarat bahwa pemenang permainan harus membayar
$b$ kepada pihak yang kalah sebagai hadiah hiburan. Bagi pemain A, nilainya tetap
$(a + b)/2$. Namun, apa hasil yang mungkin baginya dalam permainan yang diubah ini?
Jika menang, ia memperoleh seluruh taruhan $a + b$ dan harus membayar $b$ kepada B,
sehingga akhirnya memperoleh $a$. Jika kalah, ia memperoleh $b$ dari pemain B. Jadi,
pemain A memenangkan $a$ atau $b$ dengan peluang yang sama, dan nilai permainan
baginya adalah $(a + b)/2$.

Huygens menggambarkan bukti ini dengan sebuah contoh. Jika Anda ditawari permainan
dengan peluang yang sama untuk memenangkan 2~atau~8, nilai harapannya adalah~5, karena
permainan ini setara dengan permainan ketika masing-masing pemain mempertaruhkan 5 dan
sepakat membayar 3 kepada pihak yang kalah --- permainan yang nilainya jelas 5.

Proposisi kedua Huygens adalah

\begin{quote} Prop.~II. Jika saya mengharapkan $a$,~$b$, atau~$c$, yang masing-masing
dapat terjadi dengan kemudahan yang sama, maka Nilai Harapan saya adalah
$(a + b + c)/3$, atau sepertiga dari Jumlah $a$,~$b$, dan~$c$.\footnote{ibid.,
p.~35.}
\end{quote}

Argumennya di sini serupa. Tiga pemain, A,~B, dan~C, masing-masing mempertaruhkan
$$(a+b+c)/3$$ 
dalam permainan yang masing-masing mempunyai peluang sama untuk dimenangkan. Nilai
permainan ini bagi pemain A jelas sebesar jumlah yang dipertaruhkannya. Lebih lanjut,
nilai ini tidak berubah jika A membuat perjanjian dengan B bahwa jika salah satu dari
mereka menang, pemenang membayar hadiah hiburan $b$ kepada yang lain, serta membuat
perjanjian serupa dengan C dengan hadiah hiburan $c$. Berdasarkan simetri,
perjanjian-perjanjian ini tidak mengubah nilai permainan. Dalam permainan yang diubah
ini, jika A menang, ia memenangkan seluruh taruhan $a + b + c$ dikurangi hadiah
hiburan $b + c$, sehingga kemenangan akhirnya adalah $a$. Jika B menang, A
memenangkan $b$, dan jika C menang, A memenangkan $c$. Jadi, A berada dalam permainan
bernilai $(a + b + c)/3$ dengan hasil $a$,~$b$, dan~$c$ yang terjadi dengan peluang
sama. Hal ini membuktikan Proposisi II.
\par
Secara lebih umum, penalaran ini menunjukkan bahwa jika terdapat $n$ hasil
$$a_1,\ a_2,\ \ldots,\ a_n\ ,$$ 
yang semuanya terjadi dengan peluang yang sama, nilai harapannya adalah
$$
\frac {a_1 + a_2 +\cdots+ a_n}n\ .
$$

Dalam proposisi ketiganya, Huygens meninjau kasus ketika Anda memenangkan $a$ atau
$b$, tetapi dengan peluang yang tidak sama. Ia menganggap terdapat $p$ kesempatan
untuk memenangkan $a$ dan $q$ kesempatan untuk memenangkan $b$, semuanya mempunyai
peluang yang sama. Kemudian ia menunjukkan bahwa nilai harapannya adalah
$$ E = \frac p{p + q} \cdot a + \frac q{p + q} \cdot b\ .
$$ Hal ini diperoleh dengan meninjau pertaruhan yang setara dengan $p + q$ hasil yang
semuanya terjadi dengan peluang sama, dengan pembayaran $a$ pada $p$ hasil dan $b$
pada $q$ hasil. Cara ini memungkinkan Huygens menghitung nilai harapan bagi percobaan
dengan peluang yang tidak sama, setidaknya ketika peluang-peluang tersebut berupa
bilangan rasional.
\par
Jadi, alih-alih mendefinisikan nilai harapan sebagai rata-rata berbobot, Huygens
menganggap nilai harapan pertaruhan simetris tertentu telah diketahui, lalu menurunkan
nilai-nilai lainnya dari sana. Walaupun memerlukan cukup banyak manipulasi cerdik,
Huygens pada akhirnya memperoleh nilai-nilai yang sesuai dengan definisi modern kita
tentang nilai harapan. Salah satu keunggulan metode ini adalah bahwa metode tersebut
memberikan pembenaran bagi nilai harapan dalam kasus ketika tidak wajar menganggap
bahwa percobaan dapat diulang sangat banyak kali, misalnya dalam taruhan bahwa
sekurang-kurangnya dua presiden meninggal pada tanggal yang sama dalam setahun.
(Kenyataannya, tiga orang demikian; semuanya penandatangan Deklarasi Kemerdekaan, dan
ketiganya meninggal pada 4~Juli.)
\par
Dalam bukunya, Huygens menghitung nilai harapan permainan dengan menggunakan teknik
yang serupa dengan teknik yang kita gunakan untuk menghitung nilai harapan roulette di
Monte Carlo. Sebagai contoh, proposisi XIV-nya berbunyi:

\begin{quote} Prop.~XIV. Jika saya bermain bergiliran dengan orang lain, menggunakan
dua Dadu, dengan Syarat bahwa jika saya melempar 7 saya menang, dan jika ia melempar 6
ia menang, dengan membiarkannya melakukan Lemparan pertama: Carilah perbandingan
Bahaya saya terhadap bahayanya.\footnote{ibid., p.~47.}
\end{quote}
\par Deskripsi modern permainan ini adalah sebagai berikut. Huygens dan lawannya
bergiliran melempar sepasang dadu. Permainan berakhir jika Huygens melempar 7 atau
lawannya melempar 6. Lawannya melempar lebih dahulu. Berapakah peluang Huygens
memenangkan permainan?
\par Untuk menyelesaikan masalah ini, Huygens memisalkan $x$ sebagai peluangnya menang
ketika lawannya melempar lebih dahulu dan $y$ sebagai peluangnya menang ketika ia
melempar lebih dahulu. Pada lemparan pertama, lawannya menang pada 5 dari 36
kemungkinan. Dengan demikian,
$$ x = \frac {31}{36} \cdot y\ .
$$ Namun, ketika Huygens melempar, ia menang pada 6 dari 36 hasil yang mungkin, dan
pada 30 hasil lainnya ia kembali ke keadaan ketika peluangnya adalah $x$. Jadi,
$$ y = \frac 6{36} + \frac {30}{36} \cdot x\ .
$$ Dari kedua persamaan ini, Huygens memperoleh $x = 31/61$.\index{HUYGENS, C.|)}

Penggunaan awal lain nilai harapan muncul dalam argumen Pascal\index{PASCAL, B.} untuk
menunjukkan bahwa orang yang rasional seharusnya percaya akan keberadaan
Tuhan.\index{keberadaan Tuhan}\footnote{Dibahas dalam I. Hacking, \emx {The Emergence
of Probability} (Cambridge: Cambridge Univ.\ Press, 1975).} Pascal mengatakan bahwa
kita harus bertaruh apakah akan percaya atau tidak percaya. Misalkan $p$ menyatakan
peluang bahwa Tuhan tidak ada. Pembahasannya menunjukkan bahwa kita memainkan
permainan dengan dua strategi, percaya dan tidak percaya, dengan pembayaran seperti
ditampilkan dalam Tabel~\ref{table 6.4}.
\begin{table}
\centering 
$$\begin{tabular}{ccc}
              & \hspace{.7in}Tuhan tidak ada & \hspace{.05in}Tuhan ada \\
              & \\
              & \hspace{.8in}$p$              & \hspace{.15in}$1 - p$ \\
\end{tabular}
$$
$$\begin{tabular}{l|c|c|}\cline{2-2} \cline{3-3} 
 percaya      & \hspace{.5in} $-u$\hspace{.5in} &\hspace{.3in} $v\hspace{.3in} $\\
\cline{2-2} \cline{3-3} 
 tidak percaya  &    0  & $-x$\\ \cline{2-2} \cline{3-3} 
\end{tabular}
$$
\caption{Pembayaran.}
\label{table 6.4}
\end{table}

Di sini $-u$ menyatakan biaya yang Anda tanggung karena melepaskan beberapa kesenangan
duniawi sebagai akibat dari kepercayaan bahwa Tuhan ada. Jika Anda tidak percaya dan
Tuhan adalah Tuhan yang membalas dendam, Anda akan kehilangan $x$. Jika Tuhan ada dan
Anda percaya, Anda akan memperoleh v. Sekarang, untuk menentukan strategi terbaik,
Anda harus membandingkan kedua nilai harapan
$$ p(-u) + (1 - p)v \qquad {\rm dan} \qquad p0 + (1 - p)(-x),
$$ lalu memilih yang lebih besar. Secara umum, pilihan tersebut bergantung pada nilai
$p$. Namun, Pascal menganggap nilai $v$ tak hingga, sehingga strategi percaya adalah
yang terbaik berapa pun peluang yang Anda tetapkan bagi keberadaan Tuhan. Contoh ini
dianggap oleh sebagian orang sebagai awal teori keputusan. Analisis keputusan semacam
ini kini muncul dalam banyak bidang dan, khususnya, merupakan bagian penting dari
diagnostik medis serta keputusan bisnis perusahaan.

Penggunaan awal lain nilai harapan adalah untuk menentukan harga
anuitas.\index{anuitas} Kajian statistika berawal dari penggunaan catatan kematian yang
disimpan di paroki-paroki London sejak 1603. Catatan tersebut memuat jumlah pembaptisan
dan pemakaman setiap minggu. Dari catatan ini, John Graunt\index{GRAUNT, J.} membuat
estimasi populasi London dan juga menyediakan data kematian pertama,\index{tabel
kematian}\footnote{ibid., p.~108.} yang ditampilkan dalam Tabel~\ref{table 6.5}.
\begin{table}
\centering
$$ 
\begin{tabular}{rr}
    & \\ \hline
 Usia &\hspace{.35in} Penyintas  \\ \hline
 0    &  100  \\
 6    &   64  \\ 16    &   40  \\ 26    &   25  \\ 36    &   16  \\ 46    &   10  \\
56    &    6  \\ 66    &    3  \\ 76    &    1  \\ \hline
\end{tabular}$$
\caption{Data kematian Graunt.}
\label{table 6.5}
\end{table}

Seperti diamati Hacking, tampaknya Graunt menyusun tabel ini dengan menganggap bahwa
setelah usia~6 terdapat peluang tetap sekitar 5/8 untuk bertahan hidup selama satu
dasawarsa lagi.\footnote{ibid., p.~109.} Sebagai contoh, dari 64~orang yang bertahan
hidup hingga usia~6, sebanyak 5/8 dari 64, atau~40, bertahan hingga~16; sebanyak 5/8
dari 40 ini, atau~25, bertahan hingga~26; dan seterusnya. Tentu saja, ia membulatkan
angkanya ke bilangan bulat orang terdekat.

Jelas, laju kematian yang konstan tidak mungkin benar di seluruh rentang tersebut, dan
tabel-tabel yang kemudian disediakan oleh Halley lebih realistis dalam hal
ini.\footnote{E.
Halley, ``An Estimate of The Degrees of Mortality of Mankind," \emx {Phil.\ Trans.\
Royal.\ Soc.,} vol.~17 (1693), pp.~596--610; 654--656.}

Suatu \emx {anuitas berjangka}\index{anuitas!berjangka} memberikan sejumlah uang tetap
selama periode $n$ tahun. Untuk menentukan harga anuitas berjangka, kita hanya perlu
mengetahui suku bunga yang sesuai. Suatu \emx {anuitas seumur hidup}\index{anuitas!seumur
hidup} memberikan jumlah tetap selama setiap tahun kehidupan pembeli. Harga yang
sesuai bagi anuitas seumur hidup adalah nilai harapan anuitas berjangka yang dievaluasi
pada masa hidup acak pembeli. Jadi, karya Huygens dalam memperkenalkan nilai harapan
serta karya Graunt dan Halley dalam menentukan tabel kematian menghasilkan metode yang
lebih rasional untuk menetapkan harga anuitas. Ini merupakan salah satu penggunaan
serius pertama teori peluang di luar rumah perjudian.

Walaupun nilai harapan kini berperan dalam setiap cabang ilmu pengetahuan, arti
pentingnya tetap bertahan di kasino. Pada 1962, buku Edward Thorp\index{THORP, E.}
\emx {Beat the Dealer}\footnote{E. Thorp, \emx {Beat the Dealer} (New York: Random
House, 1962).} memberi pembaca strategi untuk memainkan permainan kasino populer
blackjack\index{blackjack} yang menjamin nilai harapan kemenangan positif bagi pemain.
Buku ini untuk selamanya mengubah keyakinan kasino bahwa mereka tidak dapat dikalahkan.

\exercises
\begin{LJSItem}

\i\label{exer 6.1.1} Sebuah kartu diambil secara acak dari setumpuk kartu yang
bernomor 2 sampai~10. Seorang pemain menang 1~dolar jika angka pada kartu itu ganjil dan
kalah 1~dolar jika angkanya genap. Berapakah nilai harapan kemenangan pemain tersebut?

\i\label{exer 6.1.2} Sebuah kartu diambil secara acak dari setumpuk kartu remi. Jika
kartu itu merah, pemain menang 1~dolar; jika hitam, pemain kalah 2~dolar.
Carilah nilai harapan permainan tersebut.

\i\label{exer 6.1.3} Di suatu kelas terdapat 20~siswa: 3 siswa bertinggi 5' 6", 5 siswa 5'8", 4
siswa 5'10", 4 siswa 6', dan 4 siswa 6' 2". Seorang siswa dipilih secara acak. Berapakah
nilai harapan tinggi badan siswa tersebut?

\i\label{exer 6.1.4} Di Las Vegas, roda roulette memiliki~0 dan~00, lalu
angka 1~sampai~36 yang ditandai pada petak-petak berukuran sama; roda diputar dan sebuah bola berhenti secara acak di
salah satu petak. Ketika seorang pemain mempertaruhkan 1~dolar pada suatu angka, ia menerima 36~dolar jika
bola berhenti pada angka tersebut, sehingga keuntungan bersihnya 35~dolar; jika tidak, ia kehilangan
taruhan satu dolarnya. Carilah nilai harapan kemenangannya.

\i\label{exer 6.1.5} Dalam versi kedua roulette di Las Vegas, seorang pemain bertaruh
pada merah atau hitam. Separuh angka dari~1 sampai~36 berwarna merah, dan separuhnya hitam. Jika seorang
pemain mempertaruhkan satu dolar pada hitam, dan bola berhenti pada angka hitam, ia memperoleh kembali
dolarnya beserta satu dolar tambahan. Jika bola berhenti pada angka merah atau pada~0 atau~00, ia
kehilangan dolarnya. Carilah nilai harapan kemenangan dari taruhan ini.

\i\label{exer 6.1.6} Sebuah dadu dilempar dua kali. Misalkan $X$ menyatakan jumlah kedua
angka yang muncul, dan $Y$ selisih kedua angka tersebut (tepatnya, angka
pada lemparan pertama dikurangi angka pada lemparan kedua). Tunjukkan bahwa $E(XY) = E(X)E(Y)$. Apakah
$X$ dan $Y$ saling bebas?

\istar\label{exer 6.1.7} Tunjukkan bahwa jika $X$ dan $Y$ adalah peubah acak yang masing-masing hanya
mengambil dua nilai, dan jika $E(XY) = E(X)E(Y)$, maka $X$ dan $Y$ saling bebas.

\i\label{exer 6.1.8} Sebuah keluarga kerajaan terus mempunyai anak sampai memperoleh seorang anak laki-laki atau sampai
mempunyai tiga anak, mana pun yang terjadi lebih dahulu. Asumsikan bahwa setiap anak adalah laki-laki dengan
probabilitas 1/2. Carilah nilai harapan banyaknya anak laki-laki dalam keluarga kerajaan ini dan
nilai harapan banyaknya anak perempuan.

\i\label{exer 6.1.9} Jika lemparan pertama dalam permainan craps bukan natural maupun
craps, pemain dapat memasang taruhan tambahan, sebesar taruhan awalnya, bahwa ia akan
mencapai poin sebelum tujuh muncul. Jika poinnya empat atau sepuluh, ia dibayar
dengan rasio $2 : 1$; jika lima atau sembilan, ia dibayar dengan rasio $3 : 2$; dan jika
enam atau delapan, ia dibayar dengan rasio $6 : 5$. Carilah nilai harapan kemenangan pemain
jika ia memasang taruhan tambahan ini ketika mendapat kesempatan.

\i\label{exer 6.1.10} Dalam Contoh~\ref{exam 6.12}, asumsikan bahwa Tn. Ace memutuskan untuk membeli saham
dan menahannya sampai harganya naik 1~dolar, kemudian menjualnya dan tidak membeli lagi. Modifikasilah
program {\bf StockSystem} untuk mencari distribusi keuntungannya menurut sistem ini
setelah periode dua puluh hari. Carilah nilai harapan keuntungan dan probabilitas bahwa ia memperoleh
keuntungan.

\i\label{exer 6.1.11} Pada 26~September 1980, \emx {New York Times} melaporkan
bahwa seorang asing misterius melangkah masuk ke sebuah kasino Las Vegas, memasang satu taruhan sebesar
777{,}000 dolar pada garis ``don't pass" di meja craps, lalu pergi membawa
lebih dari 1.5 juta dolar. Dalam taruhan ``don't pass", petaruh pada dasarnya
bertaruh bersama pihak kasino. Pengecualian terjadi jika pelempar memperoleh~12 pada lemparan
pertama. Dalam hal ini, pelempar kalah dan petaruh ``don't pass" hanya menerima kembali
uang yang dipertaruhkan, alih-alih menang. Tunjukkan bahwa petaruh ``don't pass" memiliki taruhan yang lebih
menguntungkan daripada pelempar.

\i\label{exer 6.1.12} Ingatlah bahwa dalam \emx {sistem penggandaan martingal}\index{sistem taruhan
martingal} (lihat Latihan~\ref{sec 1.1}.\ref{exer 1.1.10}), pemain menggandakan
taruhannya setiap kali kalah. Misalkan Anda
bermain roulette di sebuah \emx {kasino adil} tanpa 0, dan setiap kali bertaruh pada merah.
Dengan demikian, setiap kali Anda menang dengan probabilitas 1/2. Asumsikan bahwa Anda memasuki kasino dengan
100~dolar, memulai dengan
taruhan 1-dolar, dan menggunakan sistem martingal. Anda berhenti segera setelah memenangi satu taruhan, atau jika peristiwa yang kecil kemungkinannya terjadi, yaitu hitam muncul enam kali berturut-turut,
sehingga Anda telah rugi 63~dolar dan tidak dapat memasang taruhan 64-dolar yang diperlukan. Carilah
nilai harapan kemenangan Anda dengan sistem permainan ini.

\i\label{exer 6.1.13} Anda memiliki 80~dolar dan memainkan permainan berikut. Sebuah kotak
berisi dua bola putih dan dua bola hitam. Anda mengambil bola satu demi satu
tanpa pengembalian sampai semua bola habis. Pada setiap pengambilan, Anda mempertaruhkan separuh
kekayaan Anda saat itu bahwa Anda akan mengambil bola putih. Berapakah nilai harapan kekayaan akhir Anda?

\i\label{exer 6.1.14} Dalam masalah penitipan topi (lihat Contoh~\ref{exam 3.13}),
diasumsikan bahwa $N$ orang menitipkan topi mereka dan topi-topi itu dikembalikan
secara acak. Misalkan $X_j = 1$ jika orang ke-$j$ menerima topinya sendiri dan 0
jika tidak. Carilah $E(X_j)$ dan $E(X_j \cdot X_k)$ untuk $j$ yang tidak sama dengan $k$. Apakah $X_j$
dan $X_k$ saling bebas?

\i\label{exer 6.1.16} Sebuah kotak berisi dua bola emas dan tiga bola perak. Anda
diperbolehkan mengambil bola secara berurutan dan acak dari kotak tersebut. Anda menang 1~dolar
setiap kali mengambil bola emas dan kalah 1~dolar setiap kali mengambil bola perak.
Setelah diambil, bola tidak dikembalikan. Tunjukkan bahwa jika Anda terus mengambil sampai unggul
1~dolar atau sampai tidak ada lagi bola emas, permainan ini menguntungkan.

\i\label{exer 6.1.17} Gerolamo Cardano\index{CARDANO, G.}, dalam bukunya \emx {The Gambling Scholar,}
yang ditulis pada awal 1500-an, membahas permainan karnaval berikut. Terdapat enam
dadu. Setiap dadu memiliki lima sisi kosong. Sisi keenam memuat sebuah angka antara
1~dan~6---angka yang berbeda pada setiap dadu. Keenam dadu dilempar dan pemain
memenangi hadiah yang bergantung pada jumlah angka yang muncul.

\begin{enumerate}
\item Carilah, seperti yang dilakukan Cardano, nilai harapan jumlah tersebut tanpa mencari distribusinya.

\item Hadiah besar diberikan untuk jumlah yang besar, sedangkan biaya memainkan permainan ini tidak mahal.
Jelaskan mengapa hal ini dapat dilakukan.
\end{enumerate}

\i\label{exer 6.1.18} Misalkan $X$ adalah saat pertama suatu \emx {kegagalan} terjadi dalam
barisan tak hingga percobaan Bernoulli dengan probabilitas keberhasilan $p$. Misalkan $p_k
= P(X = k)$ untuk $k = 1$,~2, \dots. Tunjukkan bahwa $p_k = p^{k - 1}q$ dengan $q = 1 - p$.
Tunjukkan bahwa $\sum_k p_k = 1$. Tunjukkan bahwa $E(X) = 1/q$. Berapakah nilai harapan banyaknya
lemparan koin yang diperlukan untuk memperoleh ekor pertama?

\i\label{exer 6.1.19} Tepat satu dari enam kunci yang serupa dapat membuka suatu pintu. Jika
Anda mencoba kunci-kunci itu satu demi satu, berapakah nilai harapan banyaknya kunci yang harus
Anda coba sebelum berhasil?

\i\label{exer 6.1.20} Suatu ujian pilihan ganda diberikan. Sebuah soal memiliki empat
kemungkinan jawaban, dan tepat satu jawaban benar. Siswa diperbolehkan
memilih suatu subhimpunan dari keempat kemungkinan jawaban sebagai jawabannya. Jika subhimpunan yang dipilih
memuat jawaban benar, siswa memperoleh tiga poin, tetapi kehilangan satu
poin untuk setiap jawaban salah dalam subhimpunan pilihannya. Tunjukkan bahwa jika ia sekadar menebak suatu
subhimpunan secara seragam dan acak, nilai harapan skornya adalah nol.

\i\label{exer 6.1.21} Anda ditawari permainan berikut: sebuah koin seimbang
dilempar sampai kepala muncul untuk pertama kali (lihat Contoh~\ref{exam 6.055}). Jika
hal ini terjadi pada lemparan pertama, Anda menerima 2~dolar; jika terjadi pada lemparan kedua,
Anda menerima $2^2 = 4$ dolar; dan secara umum, jika kepala muncul untuk pertama kali
pada lemparan ke-$n$, Anda menerima $2^n$ dolar.

\begin{enumerate}
\item Tunjukkan bahwa nilai harapan kemenangan Anda tidak ada (yakni, diberikan oleh
suatu jumlah divergen) untuk permainan ini. Apakah ini berarti bahwa permainan ini menguntungkan berapa pun
biaya yang Anda bayarkan untuk memainkannya?

\item Asumsikan bahwa Anda hanya menerima $2^{10}$ dolar jika banyaknya lemparan yang diperlukan untuk memperoleh kepala pertama
lebih besar dari atau sama dengan sepuluh. Tunjukkan bahwa nilai harapan Anda
untuk permainan yang dimodifikasi ini berhingga dan carilah nilainya.

\item Asumsikan bahwa Anda membayar 10~dolar untuk setiap kali memainkan permainan asli. Tulislah sebuah
program untuk menyimulasikan 100 kali permainan dan lihatlah hasil Anda.

\item Sekarang asumsikan bahwa utilitas dari $n$ dolar adalah $\sqrt n$. Tulislah suatu bentuk
untuk nilai harapan utilitas pembayaran tersebut, dan tunjukkan bahwa bentuk ini memiliki nilai
berhingga. Perkirakan nilai ini. Ulangi latihan ini untuk kasus ketika fungsi
utilitasnya adalah $\log(n)$.

\end{enumerate}

\i\label{exer 6.1.100} Misalkan $X$ adalah peubah acak yang
berdistribusi Poisson dengan parameter $\lambda$. Tunjukkan bahwa $E(X) = \lambda$. \emx {Petunjuk}:
Ingat bahwa
$$e^x = 1 + x + \frac{x^2}{2!} + \frac{x^3}{3!} + \cdots\,.$$

\i\label{exer 6.1.22} Ingatlah bahwa dalam Latihan~\ref{sec 1.1}.\ref{exer 1.1.14}, kita
meninjau sebuah kota dengan dua rumah sakit.\index{rumah sakit} Di rumah sakit besar, sekitar 45 bayi
lahir setiap hari, sedangkan di rumah sakit kecil sekitar 15 bayi\index{bayi} lahir setiap hari. Kita
tertarik untuk menebak rumah sakit mana yang secara rata-rata memiliki jumlah hari terbanyak
dengan sifat bahwa lebih dari 60 persen anak yang lahir pada hari tersebut adalah
laki-laki. Untuk setiap rumah sakit, carilah nilai harapan banyaknya hari dalam setahun yang memiliki
sifat bahwa lebih dari 60~persen anak yang lahir pada hari tersebut adalah laki-laki.

\i\label{exer 6.1.23} Sebuah perusahaan asuransi memiliki 1{,}000 polis untuk pria berusia~50.
Perusahaan tersebut memperkirakan bahwa probabilitas seorang pria berusia~50 meninggal dalam waktu satu tahun adalah
.01. Perkirakan banyaknya klaim yang dapat diharapkan perusahaan dari para penerima manfaat pria-pria
tersebut dalam waktu satu tahun.

\i\label{exer 6.1.24} Dengan menggunakan tabel kehidupan tahun 1981 dalam Lampiran~C, tulislah
sebuah program untuk menghitung nilai harapan lama hidup pria dan wanita pada setiap kemungkinan usia
dari~1 sampai~85. Bandingkan hasil untuk pria dan wanita. Berikan komentar apakah premi asuransi
jiwa seharusnya berbeda bagi pria dan wanita.

\istar\label{exer 6.1.25} Satu set kartu ESP\index{ESP} terdiri dari 20 kartu dalam dua jenis:
misalnya sepuluh bintang dan sepuluh lingkaran (biasanya terdapat lima jenis). Set kartu dikocok dan
kartu dibuka satu demi satu. Anda, orang yang dianggap memiliki kemampuan persepsi, harus menyebutkan
simbol pada setiap kartu \emx {sebelum} kartu itu dibuka.
\par
Misalkan bahwa Anda sebenarnya hanya menebak kartu-kartu tersebut. Jika Anda tidak dapat melihat
setiap kartu setelah membuat tebakan, nilai harapan banyaknya tebakan benar mudah
dihitung, yaitu sepuluh.
\par
Sebaliknya, jika Anda menebak dengan informasi, yaitu jika Anda melihat setiap
kartu setelah menebak, tentu saja Anda mungkin berharap memperoleh skor yang lebih tinggi. Hal ini
memang benar, tetapi menghitung nilai harapan yang tepat tidak lagi mudah.
\par
Namun, simulasi komputer atas penebakan dengan informasi ini mudah dilakukan, sehingga kita
dapat memperoleh perkiraan yang baik atas nilai harapannya melalui simulasi. (Ini serupa dengan cara
pemain blackjack terampil menjadikan blackjack sebagai permainan yang menguntungkan dengan mengamati
kartu-kartu yang telah dimainkan. Lihat Latihan~\ref{exer 6.1.29}.)

\begin{enumerate}
\item Pertama, lakukan simulasi penebakan tanpa informasi dengan mengulangi
percobaan sekurang-kurangnya 1000 kali. Perkirakan nilai harapan banyaknya jawaban benar dan
bandingkan hasil Anda dengan nilai harapan teoretis.

\item Apakah strategi terbaik untuk menebak dengan informasi?

\item Lakukan simulasi penebakan dengan informasi menggunakan strategi dalam (b).
Ulangi percobaan sekurang-kurangnya 1000 kali, dan perkirakan nilai harapan dalam kasus ini.

\item Misalkan $S$ adalah banyaknya bintang dan $C$ banyaknya lingkaran dalam set kartu. Misalkan
$h(S,C)$ adalah nilai harapan kemenangan dengan menggunakan strategi penebakan optimal dalam (b). Tunjukkan
bahwa $h(S,C)$ memenuhi relasi rekursi
$$ h(S,C) = \frac S{S + C} h(S - 1,C) + \frac C{S + C} h(S,C - 1) + \frac
{\max(S,C)}{S + C}\ ,
$$ dan $h(0,0) = h(-1,0) = h(0,-1) = 0$. Dengan menggunakan relasi ini, tulislah sebuah program untuk
menghitung $h(S,C)$ dan mencari $h(10,10)$. Bandingkan nilai $h(10,10)$ yang dihitung dengan
hasil simulasi Anda dalam (c). Untuk pembahasan lebih lanjut tentang latihan ini dan
Latihan~\ref{exer 6.1.26}, lihat Diaconis dan Graham.\index{DIACONIS, P.}\index{GRAHAM,
R.}\footnote{P. Diaconis and R. Graham, ``The Analysis of Sequential Experiments with Feedback to
Subjects," \emx {Annals of Statistics,} vol.~9 (1981), pp.~3--23.}
\end{enumerate}

\istar\label{exer 6.1.26} Tinjaulah masalah ESP\index{ESP} seperti yang dijelaskan dalam Latihan~\ref{exer
6.1.25}. Sekali lagi Anda menebak dengan informasi, dan Anda menggunakan strategi
penebakan optimal: tebak \emx {bintang} jika set kartu yang tersisa memuat lebih banyak bintang, {\em
lingkaran} jika lebih banyak lingkaran, dan lempar koin jika banyaknya bintang dan lingkaran
sama. Asumsikan bahwa $S \geq C$, dengan $S$ menyatakan banyaknya bintang dan $C$ banyaknya
lingkaran.
\par Kita dapat menggambarkan hasil suatu permainan biasa pada sebuah grafik, dengan sumbu horizontal
menyatakan banyaknya langkah dan sumbu vertikal menyatakan \emx {selisih}
antara banyaknya bintang dan banyaknya lingkaran yang telah dibuka. Sebuah
permainan biasa ditunjukkan dalam Gambar~\ref{fig 6.4}. Dalam permainan khusus ini, urutan
kartu yang dibuka adalah $(C,S,S,S,S,C,C,S,S,C)$. Jadi, dalam permainan khusus
ini, terdapat enam bintang dan empat lingkaran dalam set kartu. Ini berarti, khususnya,
bahwa setiap permainan dengan set kartu ini akan memiliki grafik yang berakhir pada titik
$(10, 2)$. Kita definisikan garis $L$ sebagai garis horizontal yang melalui
titik akhir pada grafik (jadi koordinat vertikalnya hanyalah selisih antara
banyaknya bintang dan lingkaran dalam set kartu).

\putfig{4truein}{PSfig6-4}{Jalan acak untuk ESP.}{fig 6.4}

\begin{enumerate}
\item Tunjukkan bahwa ketika jalan acak berada di bawah garis $L$, pemain menebak dengan benar
saat grafik naik (bintang dibuka), dan ketika jalan berada di atas garis tersebut,
pemain menebak dengan benar saat jalan turun (lingkaran dibuka). Dari sifat ini,
tunjukkan bahwa peserta pasti memiliki sekurang-kurangnya $S$ tebakan benar.

\item Ketika jalan berada pada titik $(x,x)$ \emx {di} garis $L$, banyaknya bintang
dan lingkaran yang tersisa sama, sehingga peserta melempar koin. Tunjukkan bahwa
probabilitas jalan tersebut mencapai $(x,x)$ adalah
$$
\frac{{S \choose x}{C \choose x}}{{{S + C} \choose {2x}}}\ .
$$ \emx {Petunjuk}: Hasil dari $2x$ kartu memiliki distribusi hipergeometrik (lihat
Bagian~\ref{sec 5.1}).

\item Dengan menggunakan hasil (a) dan (b), tunjukkan bahwa nilai harapan banyaknya
tebakan benar dengan penebakan cerdas adalah
$$ S + \sum_{x = 1}^C \frac12 \frac{{S \choose x}{C \choose x}}{{{S + C} \choose
{2x}}}\ .
$$
\end{enumerate}

\i\label{exer 6.1.27} Konon\footnote{J.~F. Box,  \emx {R.~A. Fisher, The
Life of a Scientist} (New York: John Wiley and Sons, 1978).}, Dr.~B.~Muriel
Bristol menolak secangkir teh\index{teh} karena ia lebih menyukai cangkir yang terlebih dahulu dituangi
susu\index{susu}. Statistikawan terkenal R.~A. Fisher\index{FISHER, R. A.}
melakukan pengujian untuk melihat apakah ia dapat membedakan susu yang dituangkan sebelum atau sesudah teh.
Asumsikan bahwa dalam pengujian tersebut Dr.~Bristol diberi delapan cangkir teh---empat dengan susu yang dituangkan
sebelum teh dan empat dengan susu yang dituangkan sesudah teh.
% (cf. Exercise~\ref{sec 3.2}.\ref{exer 3.2.8}).

\begin{enumerate}
\item Berapakah nilai harapan banyaknya tebakan benar yang dibuat wanita tersebut jika ia tidak memperoleh
informasi setelah setiap pengujian dan hanya menebak?

\item Dengan menggunakan hasil Latihan~\ref{exer 6.1.26}, carilah nilai harapan banyaknya
tebakan benar jika ia diberi tahu hasil setiap tebakan dan menggunakan strategi penebakan
optimal.
\end{enumerate}

\i\label{exer 6.1.28} Dalam sebuah permainan komputer populer, komputer memilih suatu bilangan bulat
dari~1 sampai~$n$ secara acak. Pemain diberi $k$ kesempatan untuk menebak angka tersebut. Setelah
setiap tebakan, komputer menjawab ``benar," ``terlalu kecil," atau ``terlalu besar."

\begin{enumerate}
\item Tunjukkan bahwa jika $n \leq 2^k - 1$, maka terdapat strategi yang menjamin Anda
akan menebak angka tersebut dengan benar dalam $k$ percobaan.

\item Tunjukkan bahwa jika $n \geq 2^k - 1$, terdapat strategi yang memastikan Anda
dapat mengidentifikasi salah satu dari $2^k - 1$ angka dan dengan demikian memberikan probabilitas kemenangan sebesar $(2^k - 1)/n$.
Mengapa strategi ini optimal? Ilustrasikan hasil Anda untuk
kasus $n = 9$ dan $k = 3$.
\end{enumerate}

\i\label{exer 6.1.29} Dalam permainan kasino blackjack\index{blackjack}, bandar dibagikan dua
kartu, satu terbuka dan satu tertutup, dan setiap pemain dibagikan dua kartu yang keduanya
tertutup. Jika kartu bandar yang tampak adalah as, pemain dapat melihat kartu tertutupnya lalu
memasang taruhan yang disebut taruhan \emx {asuransi}. (Pemain ahli akan memahami mengapa taruhan ini
disebut asuransi.) Jika memasang taruhan ini, Anda akan memenanginya jika kartu kedua
bandar adalah \emx {kartu bernilai sepuluh}: yaitu sepuluh, jack, ratu, atau raja. Jika menang, Anda
dibayar dua kali taruhan asuransi; jika tidak, Anda kehilangan taruhan ini. Tunjukkan bahwa jika satu-satunya
kartu yang dapat Anda lihat adalah as milik bandar dan dua kartu Anda, dan jika kartu Anda bukan
kartu bernilai sepuluh, maka taruhan asuransi tidak menguntungkan. Namun, tunjukkan bahwa jika Anda
memainkan dua tangan sekaligus dan tidak memiliki kartu bernilai sepuluh, taruhan tersebut
menguntungkan. (Thorp\index{THORP, E.}\footnote{E.~Thorp,  \emx {Beat the Dealer} (New York: Random
House, 1962).} telah menunjukkan bahwa permainan blackjack menguntungkan bagi pemain jika ia dapat
mengingat dengan cukup baik kartu-kartu yang telah dimainkan.)

\i\label{exer 6.1.30} Asumsikan bahwa setiap kali membeli sekotak Wheaties\index{Wheaties}, Anda
menerima gambar salah satu dari $n$ pemain New York Yankees\index{New York Yankees} (lihat
Latihan~\ref{sec 3.2}.\ref{exer 3.2.34}). Misalkan $X_k$ adalah
banyaknya kotak tambahan yang harus Anda beli setelah memperoleh $k - 1$ gambar berbeda,
untuk memperoleh gambar baru berikutnya. Jadi $X_1 = 1$, sedangkan $X_2$ adalah banyaknya
kotak yang dibeli setelah itu untuk memperoleh gambar yang berbeda dari gambar pertama yang
diperoleh, dan seterusnya.

\begin{enumerate}
\item Tunjukkan bahwa $X_k$ memiliki distribusi geometrik dengan $p = (n - k + 1)/n$.

\item Simulasikan percobaan untuk tim dengan 26 pemain (25 akan lebih tepat,
tetapi kita menginginkan bilangan genap). Lakukan sejumlah simulasi dan perkirakan
nilai harapan waktu yang diperlukan untuk memperoleh 13 pemain pertama serta nilai harapan waktu untuk memperoleh
13 pemain kedua. Bagaimana perbandingan kedua nilai harapan ini?

\item Tunjukkan bahwa jika terdapat $2n$ pemain, nilai harapan waktu untuk memperoleh separuh pertama
pemain adalah
$$ 2n \left( \frac 1{2n} + \frac 1{2n - 1} +\cdots+ \frac 1{n + 1} \right)\ ,
$$ dan nilai harapan waktu untuk memperoleh separuh kedua adalah
$$ 2n \left( \frac 1n + \frac 1{n - 1} +\cdots+ 1 \right)\ .
$$

\item Dalam Contoh~\ref{exam 6.5}, kita menyatakan bahwa
$$ 1 + \frac 12 + \frac 13 +\cdots+ \frac 1n \sim \log n + .5772 + \frac 1{2n}\ .
$$ Gunakan ini untuk memperkirakan bentuk dalam (c). Bandingkan perkiraan tersebut dengan
nilai eksaknya dan juga dengan perkiraan yang Anda peroleh melalui simulasi untuk kasus $n =
26$.
\end{enumerate}

\istar\label{exer 6.1.31} (Feller\index{FELLER, W.}\footnote{W. Feller,  \emx {Introduction to
Probability Theory and Its Applications,} 3rd ed., vol.~1 (New York: John Wiley and
Sons, 1968), p.~240.}) Sejumlah besar orang, $N$, menjalani tes\index{tes darah}
darah. Tes ini dapat diberikan dengan dua cara: (1)~Setiap orang dapat diuji secara terpisah; dalam
hal ini diperlukan $N$ tes, (2)~sampel darah dari $k$ orang dapat digabungkan
dan dianalisis bersama. Jika tes ini \emx {negatif,} satu tes tersebut cukup untuk
$k$ orang. Jika tesnya \emx {positif,} masing-masing dari $k$ orang harus
diuji secara terpisah, sehingga seluruhnya diperlukan $k + 1$ tes untuk $k$ orang tersebut. Asumsikan
bahwa probabilitas $p$ suatu tes memberikan hasil positif sama untuk semua orang dan bahwa
peristiwa-peristiwa ini saling bebas.

\begin{enumerate}
\item Carilah probabilitas bahwa tes atas sampel gabungan dari $k$ orang akan memberikan hasil
positif.

\item Berapakah nilai harapan banyaknya tes $X$ yang diperlukan menurut
rencana~(2)? (Asumsikan bahwa $N$ habis dibagi $k$.)

\item Untuk $p$ yang kecil, tunjukkan bahwa nilai $k$ yang meminimumkan nilai harapan
banyaknya tes menurut rencana kedua kira-kira $1/\sqrt p$.
\end{enumerate}

\i\label{exer 6.1.32} Tulislah sebuah program untuk menjumlahkan bilangan-bilangan acak yang dipilih dari $[0,1]$
sampai jumlahnya untuk pertama kali lebih besar dari satu. Mintalah program Anda mengulangi
percobaan ini beberapa kali untuk memperkirakan nilai harapan banyaknya pemilihan yang diperlukan
agar jumlah bilangan yang dipilih untuk pertama kali melebihi~1. Berdasarkan
percobaan Anda, berapakah perkiraan Anda untuk bilangan ini?

\istar\label{exer 6.1.33} Masalah diskret terkait berikut juga memberikan petunjuk yang baik
untuk jawaban Latihan~\ref{exer 6.1.32}. Pilihlah secara acak dengan pengembalian
$t_1$,~$t_2$,
\dots,~$t_r$ dari himpunan $(1/n, 2/n, \dots, n/n)$. Misalkan $X$ adalah nilai terkecil
$r$ yang memenuhi
$$ t_1 + t_2 +\cdots+ t_r > 1\ .
$$ Maka $E(X) = (1 + 1/n)^n$. Untuk membuktikan ini, kita dapat pula memilih
$t_1$,~$t_2$, \dots,~$t_r$ secara acak dengan pengembalian dari himpunan $(1, 2,
\dots, n)$ dan memisalkan $X$ sebagai nilai terkecil $r$ sehingga
$$ t_1 + t_2 +\cdots+ t_r > n\ .
$$
\begin{enumerate}
\item Gunakan Latihan~\ref{sec 3.2}.\ref{exer 3.2.35.5} untuk menunjukkan bahwa
$$ P(X \geq j + 1) = {n \choose j}{\Bigl(\frac {1}{n}\Bigr)^j}\ .
$$
\item Tunjukkan bahwa
$$ E(X) = \sum_{j = 0}^n P(X \geq j + 1)\ .
$$
\item Dari kedua fakta ini, carilah suatu bentuk untuk $E(X)$. Bukti ini dikemukakan oleh
Harris Schultz.\index{SCHULTZ, H.}\footnote{H. Schultz, ``An Expected Value Problem," \emx {Two-Year
Mathematics Journal,} vol.~10, no.~4 (1979), pp.~277--78.}
\end{enumerate}

\istar\label{exer 6.1.34} (Kotak Korek Api Banach\index{Kotak Korek Api Banach}\footnote{W.~Feller, {\em
Introduction to Probability Theory,} vol. 1, p.~166.}) Seorang pria membawa sekotak korek api yang mula-mula berisi $N$ batang di masing-masing dari kedua
saku depannya. Setiap kali memerlukan korek api, ia memilih
salah satu saku secara acak dan mengambil sebatang dari kotak tersebut. Pada suatu hari ia merogoh salah satu saku
dan mendapati kotaknya kosong.

\begin{enumerate}
\item Misalkan $p_r$ menyatakan probabilitas bahwa saku lainnya berisi $r$ batang korek api.
Definisikan suatu barisan peubah acak \emx {pencacah} sebagai berikut: Misalkan
$X_i = 1$ jika pengambilan ke-$i$ berasal dari saku kiri, dan 0 jika berasal dari saku kanan.
Tafsirkan $p_r$ dalam $S_n = X_1 + X_2 +\cdots+ X_n$. Carilah suatu bentuk
binomial untuk $p_r$.

\item Tulislah program komputer untuk menghitung $p_r$, serta probabilitas bahwa
saku lainnya memuat sekurang-kurangnya $r$ batang korek api, untuk $N = 100$ dan $r$ dari~0 sampai~50.

\item Tunjukkan bahwa $(N - r)p_r = (1/2)(2N + 1)p_{r + 1} - (1/2)(r + 1)p_{r + 1}$\ .

\item Hitunglah $\sum_r p_r$.

\item Gunakan (c) dan (d) untuk menentukan nilai harapan $E$ dari distribusi
$\{p_r\}$.

\item Gunakan rumus Stirling untuk memperoleh pendekatan bagi $E$. Berapa banyak batang korek api
yang harus dimuat setiap kotak agar nilai harapan $E$ bernilai sekitar 13?
(Ambil $\pi = 22/7$.)
\end{enumerate}

\i\label{exer 6.1.35} Sebuah koin dilempar sampai kepala muncul untuk pertama kali. Jika
hal ini terjadi pada lemparan ke-$n$ dan $n$ ganjil, Anda menang $2^n/n$, tetapi jika $n$ genap,
Anda kalah
$2^n/n$. Maka, jika nilai harapan kemenangan Anda ada, nilainya diberikan oleh deret konvergen
$$ 1 - \frac 12 + \frac 13 - \frac 14 +\cdots
$$ yang disebut \emx {deret harmonik berselang-seling.} Mungkin kita tergoda untuk mengatakan bahwa nilai ini
seharusnya merupakan nilai harapan percobaan. Tunjukkan bahwa jika kita melakukannya,
nilai harapan suatu percobaan akan bergantung pada urutan hasil-hasilnya
didaftarkan.

\i\label{exer 6.1.37} Misalkan kita memiliki sebuah kotak yang berisi $c$
bola kuning dan $d$ bola hijau. Kita mengambil $k$ bola tanpa pengembalian
dari kotak tersebut. Carilah nilai harapan banyaknya bola kuning yang terambil. \emx {Petunjuk}: Tuliskan
banyaknya bola kuning yang terambil sebagai jumlah dari $c$ peubah acak.

\i\label{exer 6.1.38} Pembaca dipersilakan melihat Contoh~\ref{exam 6.7} untuk penjelasan
tentang berbagai opsi yang tersedia dalam roulette Monte Carlo.
\begin{enumerate}
\item Hitunglah nilai harapan kemenangan dari taruhan 1 franc pada merah menurut opsi (a).
\item Ulangi bagian (a) untuk opsi (b).
\item Bandingkan nilai harapan kemenangan untuk ketiga opsi.
\end{enumerate}

\istar\label{exer 6.1.39}
(dari Pittel\footnote{B. Pittel, Problem \#1195, \emx {Mathematics Magazine,} vol.~58,
no.\ 3 (May 1985), pg.~183.})\index{PITTEL, B.} Sebanyak $n$ buku telepon\index{buku telepon}
disimpan dalam suatu tumpukan. Probabilitas bahwa buku bernomor $i$ (dengan $1 \le i \le
n$) digunakan untuk suatu panggilan telepon adalah $p_i > 0$, dengan jumlah seluruh $p_i$ sama dengan 1. Setelah sebuah
buku digunakan, buku tersebut diletakkan di puncak tumpukan. Asumsikan bahwa panggilan-panggilan itu saling bebas,
berjarak waktu sama, dan sistem ini telah digunakan sejak waktu yang tak berhingga jauhnya di masa lalu. Misalkan
$d_i$ adalah kedalaman rata-rata buku $i$ dalam tumpukan. Tunjukkan bahwa $d_i \le d_j$ setiap kali $p_i
\ge p_j$. Jadi, secara rata-rata, buku yang lebih populer cenderung berada lebih dekat ke
puncak tumpukan. \emx {Petunjuk}: Misalkan $p_{ij}$ menyatakan probabilitas bahwa buku $i$ berada di atas
buku $j$. Tunjukkan bahwa $p_{ij} = p_{ij}(1 - p_j) + p_{ji}p_i$.

\istar\label{exer 6.1.40}
(dari Propp\footnote{J. Propp, Problem \#1159, \emx {Mathematics Magazine} vol.~57, no.\ 1
(Feb. 1984), pg.~50.})\index{PROPP, J.} Dalam soal sebelumnya, misalkan $P$ adalah probabilitas
bahwa pada saat ini setiap buku berada di tempat yang tepat, yakni buku $i$ berada pada urutan ke-$i$ dari
puncak. Carilah rumus untuk $P$ dalam $p_i$. Selain itu, carilah batas atas
terkecil untuk $P$ jika $p_i$ diperbolehkan berubah-ubah. \emx {Petunjuk}: Pertama, carilah probabilitas
bahwa buku 1 berada di tempat yang tepat. Kemudian carilah probabilitas bahwa buku 2 berada di tempat yang tepat
dengan syarat buku 1 berada di tempat yang tepat. Lanjutkan.

\istar\label{exer 6.1.41}
(dari H. Shultz dan B. Leonard\footnote{H. Shultz and B. Leonard, ``Unexpected Occurrences of the
Number $e$,"
\emx {Mathematics Magazine} vol.~62, no. 4 (October, 1989), pp.~269-271.})\index{SHULTZ,
H.}\index{LEONARD, B.} Suatu barisan bilangan acak dalam $[0, 1)$ dibangkitkan sampai barisan tersebut
tidak lagi naik monoton. Bilangan-bilangannya dipilih menurut distribusi seragam. Berapakah
nilai harapan panjang barisan tersebut? (Dalam menghitung panjangnya, suku yang merusak
kemonotonan ikut dihitung.) \emx {Petunjuk}: Misalkan $a_1,\ a_2,\ \ldots$ adalah barisan tersebut dan $X$ menyatakan
panjang barisan. Maka
$$P(X > k) = P(a_1 < a_2 < \cdots < a_k)\ ,$$
dan probabilitas pada ruas kanan mudah dihitung. Lebih lanjut, dapat ditunjukkan bahwa
$$E(X) = 1 + P(X > 1) + P(X > 2) + \cdots\ .$$

\i\label{exer 6.1.42}
Misalkan $T$ adalah peubah acak yang menghitung banyaknya 2-unshuffle yang dilakukan pada setumpuk $n$ kartu
sampai semua label pada kartu berbeda. Peubah acak ini telah dibahas dalam
Bagian~\ref{sec 3.3}. Dengan menggunakan Persamaan~\ref{eq 3.3.1} dalam bagian tersebut, bersama rumus
$$
E(T) = \sum_{s = 0}^\infty P(T > s)
$$
yang dibuktikan dalam Latihan~\ref{exer 6.1.33}, tunjukkan bahwa
$$
E(T) = \sum_{s = 0}^\infty \left(1 - {{2^s}\choose n}\frac {n!}{2^{sn}}\right)\ .
$$
Tunjukkan bahwa untuk $n = 52$, bentuk ini kira-kira sama dengan 11.7. (Seperti dinyatakan dalam
Bab~\ref{chp 3}, ini berarti bahwa secara rata-rata, hampir 12 kocokan riffle pada setumpuk 52 kartu
diperlukan agar prosesnya dapat dianggap acak.)


\end{LJSItem}

\section{Varians Peubah Acak Diskret}\label{sec 6.2}

Nilai harapan semakin berguna sebagai prediksi hasil suatu percobaan ketika hasil
tersebut kecil kemungkinannya menyimpang terlalu jauh dari nilai harapan. Dalam bagian
ini kita akan memperkenalkan suatu ukuran penyimpangan tersebut, yang disebut
varians.

\subsection*{Varians}

%**********Some comments we made on pg. 238 (prompted by Mark) have not been acted
%upon yet.

\begin{definition}\label{def 6.4} Misalkan $X$ adalah peubah acak bernilai numerik
dengan nilai harapan $\mu = E(X)$. Maka \emx {varians}\index{varians} dari $X$, yang dilambangkan dengan
$V(X)$, adalah
$$ V(X) = E((X - \mu)^2)\ .
$$
\end{definition} Perhatikan bahwa menurut Teorema~\ref{thm 6.3.5}, $V(X)$ diberikan oleh

\begin{equation} V(X) = \sum_x (x - \mu)^2 m(x)\ , 
\label{eq 6.1}
\end{equation} dengan $m$ sebagai fungsi distribusi dari $X$.

\subsection*{Simpangan Baku}

\emx {Simpangan baku}\index{simpangan baku} dari $X$, yang dilambangkan dengan $D(X)$, adalah $D(X) =
\sqrt {V(X)}$. Kita sering menulis $\sigma$ untuk $D(X)$ dan $\sigma^2$ untuk $V(X)$.


\begin{example}\label{exam 6.13} Tinjaulah satu kali lemparan sebuah dadu. Misalkan $X$ adalah angka
yang muncul. Untuk mencari
$V(X)$, pertama-tama kita harus mencari nilai harapan $X$. Nilainya adalah
\begin{eqnarray*}
\mu & = & E(X) = 1\Bigl(\frac 16\Bigr) + 2\Bigl(\frac 16\Bigr) + 3\Bigl(\frac
16\Bigr) +  4\Bigl(\frac 16\Bigr) + 5\Bigl(\frac 16\Bigr) + 6\Bigl(\frac 16\Bigr) \\
    & = & \frac 72\ .
\end{eqnarray*}

Untuk mencari varians $X$, kita membentuk peubah acak baru $(X - \mu)^2$ dan
menghitung nilai harapannya. Kita dapat melakukannya dengan mudah menggunakan tabel berikut.

\begin{table}
\centering
$$\begin{tabular}{ccc}
  && \\ \hline
$x$ & $m(x)$ & $(x - 7/2)^2$ \\ \hline 1 & 1/6 &       25/4  \\ 2 & 1/6
&\hspace{.05in}       9/4  \\ 3 & 1/6 &\hspace{.05in}       1/4  \\ 4 & 1/6
&\hspace{.05in}       1/4  \\ 5 & 1/6 &\hspace{.05in}       9/4  \\ 6 & 1/6 &      
25/4  \\\hline
\end{tabular}
$$ 
\caption{Perhitungan varians.}
\label{table 6.6}
\end{table}

Dari tabel ini kita memperoleh $E((X - \mu)^2)$ sebagai
\begin{eqnarray*} V(X) & = & \frac 16 \left( \frac {25}4 + \frac 94 + \frac 14 +
\frac 14 + \frac 94 + \frac {25}4 \right) \\
     & = &\frac {35}{12}\ ,
\end{eqnarray*} dan simpangan baku $D(X) = \sqrt{35/12} \approx 1.707$.
\end{example}

\subsection*{Perhitungan Varians}\index{varians!perhitungan}

Selanjutnya kita membuktikan sebuah teorema yang memberikan bentuk alternatif yang berguna untuk menghitung
varians.

\begin{theorem}\label{thm 6.7} Jika $X$ adalah sebarang peubah acak dengan $E(X) = \mu$, maka
$$ V(X) = E(X^2) - \mu^2\ .
$$
\proof Kita mempunyai
\begin{eqnarray*} V(X) & = & E((X - \mu)^2) = E(X^2 - 2\mu X + \mu^2) \\
     & = & E(X^2) - 2\mu E(X) + \mu^2 = E(X^2) - \mu^2\ .
\end{eqnarray*}
\end{theorem}

Dengan menggunakan Teorema~\ref{thm 6.7}, kita dapat menghitung varians hasil satu lemparan
dadu dengan terlebih dahulu menghitung
\begin{eqnarray*} E(X^2) & = & 1\Bigl(\frac 16\Bigr) + 4\Bigl(\frac 16\Bigr) +
9\Bigl(\frac 16\Bigr) + 16\Bigl(\frac 16\Bigr) + 25\Bigl(\frac 16\Bigr) +
36\Bigl(\frac 16\Bigr) \\
       & = &\frac {91}6\ ,
\end{eqnarray*} dan,
$$ 
V(X)  =  E(X^2) - \mu^2 = \frac {91}{6} - \Bigl(\frac 72\Bigr)^2
 = \frac {35}{12}\ ,
$$
yang sesuai dengan nilai yang diperoleh langsung dari definisi
$V(X)$.


\subsection*{Sifat-Sifat Varians}

Varians memiliki sifat-sifat yang sangat berbeda dari nilai harapan. Jika
$c$ adalah sebarang konstanta, $E(cX) = cE(X)$ dan $E(X + c) = E(X) + c$. Kedua pernyataan ini
menyiratkan bahwa nilai harapan merupakan fungsi linear. Namun, varians tidak
linear, seperti terlihat dalam teorema berikut.

\begin{theorem}\label{thm 6.6} Jika $X$ adalah sebarang peubah acak dan $c$ adalah sebarang konstanta,
maka
$$ V(cX)  = c^2 V(X)
$$ dan
$$ V(X + c) = V(X)\ .
$$
\proof Misalkan $\mu = E(X)$. Maka $E(cX) = c\mu$, dan
\begin{eqnarray*} V(cX) &=& E((cX - c\mu)^2) = E(c^2(X - \mu)^2) \\
      &=& c^2 E((X - \mu)^2) = c^2 V(X)\ .
\end{eqnarray*}

Untuk membuktikan pernyataan kedua, perhatikan bahwa untuk menghitung $V(X + c)$, kita mengganti
$x$ dengan $x + c$ dan $\mu$ dengan $\mu + c$ dalam Persamaan~\ref{eq 6.1}. Kemudian suku-suku
$c$ saling menghilangkan, sehingga tersisa $V(X)$.
\end{theorem}

Sekarang kita beralih ke beberapa sifat umum varians. Ingatlah bahwa jika $X$ dan
$Y$ adalah dua peubah acak sebarang, $E(X + Y) = E(X) + E(Y)$. Hal ini tidak selalu berlaku
untuk varians. Sebagai contoh, misalkan $X$ adalah peubah acak dengan
$V(X) \ne 0$, dan definisikan $Y = -X$. Maka $V(X) = V(Y)$, sehingga $V(X) + V(Y) =
2V(X)$. Namun,
$X + Y$ selalu 0 dan karenanya memiliki varians 0. Jadi, $V(X + Y) \ne V(X) + V(Y)$.
\par Namun, dalam kasus penting peubah-peubah acak yang saling bebas, \emx {varians
jumlah sama dengan jumlah varians.}

\begin{theorem}\label{thm 6.8} Misalkan $X$ dan $Y$ adalah dua peubah acak yang \emx {saling
bebas}. Maka
$$V(X + Y) = V(X) + V(Y)\ .$$
\proof Misalkan $E(X) = a$ dan $E(Y) = b$. Maka
\begin{eqnarray*} V(X + Y) & = & E((X + Y)^2) - (a + b)^2 \\
         & = & E(X^2) + 2E(XY) + E(Y^2) - a^2 - 2ab - b^2\ .
\end{eqnarray*} Karena $X$ dan $Y$ saling bebas, $E(XY) = E(X)E(Y) = ab$. Jadi,
$$ V(X + Y) = E(X^2) - a^2 + E(Y^2) - b^2 = V(X) + V(Y)\ .
$$
\end{theorem}

Bukti ini mudah diperluas melalui induksi matematika untuk menunjukkan bahwa {\em
varians jumlah sebarang banyaknya peubah acak yang saling bebas sama dengan jumlah
varians masing-masing.} Dengan demikian, kita memperoleh teorema berikut.

\begin{theorem}\label{thm 6.9} Misalkan $X_1$,~$X_2$, \dots,~$X_n$ merupakan suatu proses
percobaan saling bebas dengan $E(X_j) =
\mu$ dan $V(X_j) = \sigma^2$. Misalkan
$$ S_n = X_1 + X_2 +\cdots+ X_n
$$ merupakan jumlahnya, dan
$$ A_n = \frac {S_n}n
$$ merupakan rataannya. Maka
\begin{eqnarray*} E(S_n) &=& n\mu\ , \\
V(S_n) &=& n\sigma^2\ , \\ 
\sigma(S_n) &=& \sigma \sqrt{n}\ , \\
E(A_n) &=& \mu\ , \\
V(A_n) &=& \frac {\sigma^2}\ , \\
\sigma(A_n) &=& \frac{\sigma}{\sqrt n}\ .
\end{eqnarray*}
\proof Karena semua peubah acak $X_j$ mempunyai nilai harapan yang sama, kita mempunyai
$$ E(S_n) = E(X_1) +\cdots+ E(X_n) = n\mu\ ,
$$ 
$$ V(S_n) = V(X_1) +\cdots+ V(X_n) = n\sigma^2\ ,
$$
dan
$$\sigma(S_n) = \sigma \sqrt{n}\ .$$

Kita telah melihat bahwa jika suatu peubah acak $X$ dengan rataan $\mu$ dan varians
$\sigma^2$ dikalikan dengan konstanta $c$, peubah acak baru tersebut memiliki nilai harapan $c\mu$ dan
varians $c^2\sigma^2$. Jadi,
$$ E(A_n) = E\left(\frac {S_n}n \right) = \frac {n\mu}n = \mu\ ,
$$ dan
$$ V(A_n) = V\left( \frac {S_n}n \right) = \frac {V(S_n)}{n^2} = \frac
{n\sigma^2}{n^2} = \frac {\sigma^2}n\ .
$$ 
Terakhir, simpangan baku $A_n$ diberikan oleh
$$\sigma(A_n) = \frac {\sigma}{\sqrt n}\ .$$
\end{theorem}

Persamaan terakhir dalam teorema di atas menyiratkan bahwa dalam suatu proses percobaan saling bebas,
jika masing-masing suku memiliki varians berhingga, maka simpangan baku
rataannya menuju 0 ketika $n \rightarrow \infty$. Karena simpangan baku memberi kita
informasi tentang sebaran distribusi di sekitar rataan, kita melihat bahwa untuk nilai $n$
yang besar, nilai $A_n$ biasanya sangat dekat dengan rataan $A_n$, yang
sama dengan $\mu$, seperti ditunjukkan di atas. Pernyataan ini dibuat tepat dalam Bab~\ref{chp 8},
dan disebut Hukum Bilangan Besar. Sebagai contoh, misalkan $X$ menyatakan hasil lemparan
sebuah dadu seimbang. Dalam Gambar~\ref{fig 6.4.5}, kita menunjukkan distribusi peubah acak
$A_n$ yang berkaitan dengan $X$, untuk $n = 10$ dan $n = 100$.

\putfig{5truein}{PSfig6-4-5}
{Distribusi empiris $A_n$.}{fig 6.4.5}

\begin{example}\label{exam 6.14} Tinjaulah $n$ kali lemparan sebuah dadu. Kita telah melihat bahwa jika
$X_j$ adalah hasil lemparan
ke-$j$, maka $E(X_j) = 7/2$ dan $V(X_j) = 35/12$. Jadi, jika $S_n$ adalah jumlah
hasil-hasil tersebut, dan $A_n = S_n/n$ adalah rataan hasilnya, kita mempunyai
$E(A_n) = 7/2$ dan $V(A_n) = (35/12)/n$. Oleh karena itu, ketika $n$ bertambah, nilai harapan
rataan tetap konstan, tetapi varians menuju~0. Jika varians
merupakan ukuran nilai harapan penyimpangan dari rataan, hal ini menunjukkan bahwa untuk
$n$ yang besar, kita dapat mengharapkan rataan sangat dekat dengan nilai harapan. Hal ini memang
benar, dan kita akan membenarkannya dalam Bab~\ref{chp 8}.
\end{example}

\subsection*{Percobaan Bernoulli}

Selanjutnya, tinjaulah proses percobaan Bernoulli umum. Seperti biasa, kita misalkan $X_j = 1$ jika
hasil ke-$j$ merupakan keberhasilan dan 0 jika merupakan kegagalan. Jika $p$ adalah probabilitas
keberhasilan, dan $q = 1 - p$, maka
\begin{eqnarray*} E(X_j) & = & 0q + 1p = p\ , \\ E(X_j^2) & = & 0^2q + 1^2p = p\ ,
\end{eqnarray*} dan
$$ V(X_j) = E(X_j^2) - (E(X_j))^2 = p - p^2 = pq\ .
$$

Jadi, untuk percobaan Bernoulli, jika $S_n = X_1 + X_2 +\cdots+ X_n$ adalah banyaknya
keberhasilan, maka $E(S_n) = np$, $V(S_n) = npq$, dan $D(S_n) = \sqrt{npq}.$ Jika
$A_n = S_n/n$ adalah proporsi keberhasilan, maka $E(A_n) = p$, $V(A_n) = pq/n$,
dan $D(A_n) = \sqrt{pq/n}$. Kita melihat bahwa nilai harapan proporsi keberhasilan tetap
$p$ dan varians menuju~0. Hal ini menunjukkan bahwa penafsiran frekuensi atas
peluang adalah tepat. Kita akan membuat pernyataan ini lebih tepat dalam Bab~\ref{chp 8}.

\begin{example}\label{exam 6.15} Misalkan $T$ menyatakan banyaknya percobaan sampai keberhasilan
pertama dalam suatu proses percobaan Bernoulli. Maka $T$ berdistribusi geometrik. Berapakah
varians $T$? Dalam Contoh~\ref{exam 5.7}, kita melihat bahwa
$$ m_T = \pmatrix{1 &  2 &    3 & \cdots \cr p & qp & q^2p & \cdots \cr}.
$$ Dalam Contoh~\ref{exam 6.8}, kita menunjukkan bahwa
$$E(T) = 1/p\ .$$  
Jadi,
$$V(T) = E(T^2) - 1/p^2\ ,$$  
sehingga kita hanya perlu mencari
\begin{eqnarray*} E(T^2) & = & 1p + 4qp + 9q^2p + \cdots \\
       & = & p(1 + 4q + 9q^2 + \cdots )\ .
\end{eqnarray*} Untuk menghitung jumlah ini, kita mulai lagi dengan
$$ 1 + x + x^2 +\cdots= \frac 1{1 - x}\ .
$$ Dengan mendiferensialkan, kita memperoleh
$$ 1 + 2x + 3x^2 +\cdots= \frac 1{(1 - x)^2}\ .
$$ Dengan mengalikan dengan $x$,
$$ x + 2x^2 + 3x^3 +\cdots= \frac x{(1 - x)^2}\ .
$$ Dengan mendiferensialkan sekali lagi, diperoleh
$$ 1 + 4x + 9x^2 +\cdots= \frac {1 + x}{(1 - x)^3}\ .
$$ Jadi,
$$ E(T^2) = p\frac {1 + q}{(1 - q)^3} = \frac {1 + q}{p^2}
$$ dan
\begin{eqnarray*} V(T) & = & E(T^2) - (E(T))^2 \\
     & = & \frac {1 + q}{p^2} - \frac 1{p^2} = \frac q{p^2}\ .
\end{eqnarray*}

Sebagai contoh, varians banyaknya lemparan koin sampai kepala pertama
muncul adalah $(1/2)/(1/2)^2 = 2$. Varians banyaknya lemparan dadu sampai
enam pertama muncul adalah $(5/6)/(1/6)^2 = 30$. Perhatikan bahwa ketika $p$ mengecil,
varians meningkat dengan cepat. Hal ini berkaitan dengan meningkatnya sebaran
distribusi geometrik ketika $p$ mengecil (seperti dicatat dalam Gambar~\ref{fig 5.4}).
\end{example}

\subsection*{Distribusi Poisson}\index{distribusi Poisson!varians}

Seperti dalam kasus nilai harapan, varians distribusi
Poisson dengan parameter $\lambda$ mudah diduga. Ingatlah bahwa varians suatu
distribusi binomial dengan parameter $n$ dan $p$ sama dengan $npq$. Ingat pula bahwa
distribusi Poisson dapat diperoleh sebagai limit distribusi binomial jika
$n$ menuju $\infty$ dan $p$ menuju 0 sedemikian sehingga hasil kalinya tetap pada
nilai $\lambda$. Dalam kasus ini, $npq = \lambda q$ mendekati $\lambda$, karena
$q$ menuju 1. Jadi, untuk distribusi Poisson dengan parameter $\lambda$, kita patut
menduga bahwa variansnya adalah $\lambda$. Pembaca diminta menunjukkan hal ini dalam
Latihan~\ref{exer 6.2.100}.

\exercises
\begin{LJSItem}

 
\i\label{exer 6.2.1} Sebuah bilangan dipilih secara acak dari himpunan
$S = \{-1,0,1\}$. Misalkan $X$ adalah bilangan yang dipilih.  Carilah nilai harapan,
varians, dan simpangan baku $X$.

\i\label{exer 6.2.2} Peubah acak $X$ memiliki distribusi
$$ p_X = \pmatrix{ 0 & 1 & 2 & 4 \cr 1/3 & 1/3 & 1/6 & 1/6 \cr}\ .
$$ Carilah nilai harapan, varians, dan simpangan baku $X$.

\i\label{exer 6.2.3} Anda memasang taruhan 1 dolar pada bilangan 17 di Las Vegas,
sedangkan teman Anda memasang taruhan 1 dolar pada warna hitam (lihat
Latihan~\ref{sec 1.1}.\ref{exer 1.1.6}~dan~\ref{sec 1.1}.\ref{exer 1.1.7}).  Misalkan
$X$ adalah kemenangan Anda dan $Y$ adalah kemenangan teman Anda.  Bandingkan
$E(X)$,~$E(Y)$ serta $V(X)$,~$V(Y)$.  Apa yang disampaikan perhitungan ini tentang
sifat kemenangan Anda jika Anda dan teman Anda membuat serangkaian taruhan, dengan
Anda selalu bertaruh pada suatu bilangan dan teman Anda bertaruh pada suatu warna?

\i\label{exer 6.2.4} $X$ adalah peubah acak dengan $E(X) = 100$ dan $V(X) = 15$.
Carilah
\begin{enumerate}
\item $E(X^2)$.

\item $E(3X + 10)$.

\item $E(-X)$.

\item $V(-X)$.

\item $D(-X)$.
\end{enumerate}

\i\label{exer 6.2.5} Dalam suatu proses manufaktur, suhu (Fahrenheit) tidak pernah
menyimpang lebih dari $2^\circ$ dari $62^\circ$.  Suhu tersebut sebenarnya merupakan
peubah acak $F$ dengan distribusi
$$ P_F = \pmatrix{ 60 & 61 & 62 & 63 & 64 \cr 1/10 & 2/10 & 4/10 & 2/10 & 1/10 \cr}\ .
$$
\begin{enumerate}
\item Carilah $E(F)$ dan $V(F)$.
\item Definisikan $T = F - 62$.  Carilah $E(T)$ dan $V(T)$, lalu bandingkan jawaban
ini dengan jawaban pada bagian (a).

\item Diputuskan bahwa pembacaan suhu akan dilaporkan pada skala Celsius, yaitu
$C = (5/9)(F - 32)$.  Sekarang berapakah nilai harapan dan varians pembacaan tersebut?
\end{enumerate}

\i\label{exer 6.2.6} Tulislah program komputer untuk menghitung rataan dan varians
suatu distribusi yang Anda tentukan sebagai data.  Gunakan program tersebut untuk
membandingkan varians kedua densitas berikut, yang sama-sama memiliki nilai harapan~0:
$$  p_X = \pmatrix{ -2 & -1 & 0 & 1 & 2 \cr 3/11 & 2/11 & 1/11 & 2/11 & 3/11 \cr}\ ;
$$
$$ p_Y = \pmatrix{ -2 & -1 & 0 & 1 & 2 \cr 1/11 & 2/11 & 5/11 & 2/11 & 1/11 \cr}\ .
$$

\i\label{exer 6.2.7} Sebuah koin dilempar tiga kali.  Misalkan $X$ adalah banyaknya
kepala yang muncul.  Carilah $V(X)$ dan $D(X)$.

\i\label{exer 6.2.8} Sebuah sampel acak berisi 2400 orang ditanya apakah mereka
mendukung usulan pemerintah untuk membangun pembangkit listrik tenaga nuklir baru.
Jika 40~persen penduduk negara tersebut mendukung usulan ini, carilah nilai harapan
dan simpangan baku bagi banyaknya orang $S_{2400}$ dalam sampel yang mendukung usulan.

\i\label{exer 6.2.9} Sebuah dadu diberi bobot sedemikian rupa sehingga peluang
kemunculan suatu sisi sebanding dengan bilangan pada sisi tersebut.  Dadu dilempar
dengan hasil~$X$.  Carilah $V(X)$ dan $D(X)$.

\i\label{exer 6.2.10} Buktikan fakta-fakta berikut tentang simpangan baku.
\begin{enumerate}
\item $D(X + c) = D(X)$.

\item $D(cX) = |c|D(X)$.
\end{enumerate}

\i\label{exer 6.2.11} Sebuah bilangan dipilih secara acak dari bilangan bulat
1,~2,~3, \dots,~$n$.  Misalkan $X$ adalah bilangan yang dipilih.  Tunjukkan bahwa
$E(X) = (n + 1)/2$ dan $V(X) = (n - 1)(n + 1)/12$.   \emx {Petunjuk}: Identitas
berikut mungkin berguna:
$$ 1^2 + 2^2 + \cdots + n^2 = \frac{(n)(n+1)(2n+1)}{6}\ .
$$

\i\label{exer 6.2.13} Misalkan $X$ adalah peubah acak dengan $\mu = E(X)$ dan
$\sigma^2 = V(X)$.  Definisikan $X^* = (X - \mu)/\sigma$.  Peubah acak $X^*$ disebut
\emx {peubah acak baku}\index{peubah acak baku} yang berkaitan dengan $X$.  Tunjukkan
bahwa peubah acak baku ini memiliki nilai harapan 0 dan varians 1.

\i\label{exer 6.2.14} Peter dan Paul memainkan Kepala atau Ekor (lihat
Contoh~\ref{exam 1.3}).  Misalkan $W_n$ adalah kemenangan Peter setelah $n$
pertandingan.  Tunjukkan bahwa $E(W_n) = 0$ dan $V(W_n) = n$.

\i\label{exer 6.2.15} Carilah nilai harapan dan varians banyaknya anak laki-laki
serta banyaknya anak perempuan dalam sebuah keluarga kerajaan yang terus memiliki
anak sampai lahir seorang anak laki-laki atau sampai terdapat tiga anak, mana pun yang
terjadi lebih dahulu.

\i\label{exer 6.2.16} Misalkan topi milik $n$ orang dikembalikan secara acak.
Misalkan $X_i = 1$ jika orang ke-$i$ memperoleh kembali topinya sendiri dan 0 jika
tidak.  Misalkan
$S_n =
\sum_{i = 1}^n X_i$.  Maka $S_n$ adalah banyaknya orang yang memperoleh kembali
topinya sendiri.  Tunjukkan bahwa
\begin{enumerate}
\item $E(X_i^2) = 1/n$.

\item $E(X_i \cdot X_j) = 1/n(n - 1)$ untuk $i \ne j$.

\item $E(S_n^2) = 2$ (gunakan (a) dan (b)).

\item $V(S_n) = 1$.
\end{enumerate}

\i\label{exer 6.2.17} Misalkan $S_n$ adalah banyaknya keberhasilan dalam $n$
percobaan independen.  Gunakan program {\bf BinomialProbabilities}
(Bagian~\ref{sec 3.2}) untuk menghitung, bagi $n$,~$p$, dan~$j$ yang diberikan, peluang
$$ P(-j\sqrt{npq} < S_n - np < j\sqrt{npq})\ .
$$
\begin{enumerate}
\item Misalkan $p = .5$, dan hitunglah peluang ini untuk $j = 1$,~2,~3 serta $n =
10$,~30,~50.  Lakukan hal yang sama untuk $p = .2$.

\item Tunjukkan bahwa \emx {peubah acak baku} $S_n^* = (S_n - np)/\sqrt{npq}$
memiliki nilai harapan 0 dan varians 1.  Apa yang disampaikan hasil Anda dari (a)
tentang besaran baku $S_n^*$ ini?
\end{enumerate}

\i\label{exer 6.2.18} Misalkan $X$ adalah hasil suatu percobaan acak dengan
$E(X) = \mu$ dan $V(X) = \sigma^2$.  Ketika $\mu$ dan $\sigma^2$ tidak diketahui,
statistikawan sering mengestimasikannya dengan mengulangi percobaan $n$ kali dan
memperoleh hasil $x_1$,~$x_2$, \dots,~$x_n$. Parameter $\mu$ diestimasi dengan
\emx {rataan sampel}\index{rataan sampel}
$$
\bar{x} = \frac 1n \sum_{i = 1}^n x_i\ ,
$$ dan $\sigma^2$ dengan \emx {varians sampel}\index{varians sampel}
$$ s^2 = \frac 1n \sum_{i = 1}^n (x_i - \bar x)^2\ .
$$ Maka $s$ adalah \emx {simpangan baku sampel.}\index{simpangan baku sampel}
Rumus-rumus ini seharusnya mengingatkan pembaca pada definisi rataan dan varians
teoretis.  (Banyak statistikawan mendefinisikan varians sampel dengan mengganti
koefisien $1/n$ menjadi $1/(n-1)$.  Jika definisi alternatif ini digunakan, nilai
harapan $s^2$ sama dengan $\sigma^2$. Lihat Latihan~\ref{exer 6.2.19}, bagian (d).)

Tulislah program komputer yang melempar dadu $n$ kali dan menghitung rataan sampel
serta varians sampel.  Ulangi percobaan ini beberapa kali untuk~$n = 10$ dan~$n =
1000$.  Seberapa baik rataan sampel dan varians sampel mengestimasi rataan sebenarnya
7/2 dan varians 35/12?

\i\label{exer 6.2.19} Tunjukkan bahwa, untuk rataan sampel $\bar x$ dan varians
sampel $s^2$ sebagaimana didefinisikan dalam Latihan~\ref{exer 6.2.18},
\begin{enumerate}
\item $E(\bar x) = \mu$.

\item $E\bigl((\bar x - \mu)^2\bigr) = \sigma^2/n$.

\item $E(s^2) = \frac {n-1}n\sigma^2$. \emx {Petunjuk}: Untuk (c), tuliskan
\begin{eqnarray*}
\sum_{i = 1}^n (x_i - \bar x)^2 & = & \sum_{i = 1}^n \bigl((x_i - \mu) -
(\bar x - \mu)\bigr)^2 \\
     & = & \sum_{i = 1}^n (x_i - \mu)^2 - 2(\bar x - \mu) \sum_{i = 1}^n (x_i -
\mu) + n(\bar x - \mu)^2 \\
     & = & \sum_{i = 1}^n (x_i - \mu)^2 - n(\bar x - \mu)^2,
\end{eqnarray*} lalu ambil nilai harapan kedua ruas, dengan menggunakan bagian (b)
jika diperlukan.

\item Tunjukkan bahwa jika dalam definisi $s^2$ pada Latihan~\ref{exer 6.2.18} kita
mengganti koefisien $1/n$ dengan koefisien $1/(n-1)$, maka $E(s^2) = \sigma^2$.
(Hal ini menunjukkan mengapa banyak statistikawan menggunakan koefisien $1/(n-1)$.
Bilangan $s^2$ digunakan untuk mengestimasi besaran $\sigma^2$ yang tidak diketahui.
Jika suatu estimator memiliki nilai rata-rata yang sama dengan besaran yang
diestimasi, estimator itu disebut \emx {tak bias}.\index{estimator tak bias}  Jadi,
pernyataan $E(s^2) = \sigma^2$ mengatakan bahwa $s^2$ merupakan estimator tak bias
bagi $\sigma^2$.)
\end{enumerate}

\i\label{exer 6.2.20} Misalkan $X$ adalah peubah acak yang mengambil nilai
$a_1$,~$a_2$, \dots, $a_r$ dengan peluang $p_1$,~$p_2$, \dots,~$p_r$ dan memenuhi
$E(X) = \mu$.  Definisikan \emx {ukuran sebaran}\index{ukuran sebaran} $X$ sebagai berikut:
$$
\bar\sigma = \sum_{i = 1}^r |a_i - \mu|p_i\ .
$$ Besaran ini, seperti simpangan baku, merupakan cara mengukur seberapa tersebar
suatu peubah acak di sekitar rataannya.  Ingat bahwa varians jumlah peubah acak yang
saling independen sama dengan jumlah varians masing-masing peubah.  Kuadrat ukuran
sebaran berhubungan dengan varians seperti ukuran sebaran berhubungan dengan simpangan
baku.  Tunjukkan dengan sebuah contoh bahwa kuadrat ukuran sebaran jumlah dua peubah
acak independen tidak harus sama dengan jumlah kuadrat ukuran sebaran masing-masing.

\i\label{exer 6.2.21} Kita memiliki dua instrumen yang mengukur jarak antara dua
titik.  Pengukuran yang diberikan oleh kedua instrumen merupakan peubah acak $X_1$
dan $X_2$ yang independen dengan $E(X_1) = E(X_2) = \mu$, dengan $\mu$ sebagai jarak
sebenarnya.  Berdasarkan pengalaman dengan instrumen ini, kita mengetahui nilai
varians $\sigma_1^2$ dan $\sigma_2^2$.  Kedua varians ini tidak harus sama.  Dari dua
pengukuran, kita mengestimasi $\mu$ dengan rata-rata berbobot $\bar \mu = wX_1 +
(1 - w)X_2$.  Di sini $w$ dipilih dalam $[0,1]$ untuk meminimumkan varians $\bar \mu$.
\begin{enumerate}
\item Berapakah $E(\bar \mu)$?

\item Bagaimanakah $w$ harus dipilih dalam $[0,1]$ agar varians $\bar \mu$ minimum?
\end{enumerate}

\i\label{exer 6.2.22} Misalkan $X$ adalah peubah acak dengan $E(X) = \mu$ dan
$V(X) = \sigma^2$.  Tunjukkan bahwa fungsi $f(x)$ yang didefinisikan oleh
$$ f(x) = \sum_\omega (X(\omega) - x)^2 p(\omega)
$$ mencapai nilai minimum ketika $x = \mu$.

\i\label{exer 6.2.23} Misalkan $X$ dan $Y$ adalah dua peubah acak yang didefinisikan
pada ruang sampel terhingga $\Omega$.  Asumsikan bahwa $X$,~$Y$, $X + Y$, dan~$X - Y$
semuanya memiliki distribusi yang sama.  Buktikan bahwa $P(X = Y = 0) = 1$.

\i\label{exer 6.2.24} Jika $X$ dan $Y$ adalah sebarang dua peubah acak, maka {\em
kovarians} $X$ dan $Y$ didefinisikan oleh {\rm Cov}$(X,Y) = E((X - E(X))(Y -
E(Y)))$.  Perhatikan bahwa {\rm Cov}$(X,X) = V(X)$.  Tunjukkan bahwa jika $X$ dan $Y$
independen, maka {\rm Cov}$(X,Y) = 0$; lalu berikan contoh yang menunjukkan bahwa
dapat berlaku {\rm Cov}$(X,Y) = 0$ sementara $X$ dan $Y$ tidak independen.

\istar\label{exer 6.2.25} Seorang profesor ingin menyusun ujian benar-salah
\index{ujian benar-salah} dengan $n$ soal.  Ia mengasumsikan bahwa soal-soal dapat
dirancang sedemikian rupa sehingga mahasiswa menjawab soal ke-$j$ dengan benar dengan
peluang $p_j$, dan bahwa jawaban atas berbagai soal dapat dipandang sebagai percobaan
independen.  Misalkan $S_n$ adalah banyaknya soal yang dijawab benar oleh seorang
mahasiswa.  Profesor ingin memilih $p_j$ sehingga $E(S_n) = .7n$ dan varians $S_n$
sebesar mungkin.  Tunjukkan bahwa untuk mencapainya, ia harus memilih $p_j = .7$
untuk semua~$j$; artinya, semua soal harus memiliki tingkat kesulitan yang sama.

\i\label{exer 6.2.26} (Lamperti\index{LAMPERTI, J.}\footnote{Komunikasi pribadi.})
Sebuah guci berisi tepat 5000 bola. Sebanyak $X$ bola berwarna putih, dengan jumlah
ini tidak diketahui, dan sisanya merah; $X$ adalah peubah acak dengan distribusi peluang pada
bilangan bulat 0,~1,~2, \dots,~5000.
\begin{enumerate}
\item Misalkan kita mengetahui bahwa $E(X) = \mu$.  Tunjukkan bahwa informasi ini
cukup untuk menghitung peluang bahwa sebuah bola yang diambil secara acak dari guci
berwarna putih.  Berapakah peluang ini?

\item Kita mengambil sebuah bola dari guci, memeriksa warnanya, mengembalikannya,
lalu mengambil bola lain.  Dalam kondisi apakah, jika ada, hasil kedua pengambilan
independen; yakni, apakah
$$
 P({{\rm putih},{\rm putih}}) =  P({{\rm putih}})^2\ ?
$$

\item Misalkan varians $X$ adalah $\sigma^2$.  Berapakah peluang mengambil dua bola
putih dalam bagian (b)?

\end{enumerate}


\i\label{exer 6.2.27} Untuk barisan percobaan Bernoulli, misalkan $X_1$ adalah
banyaknya percobaan hingga keberhasilan pertama.  Untuk~$j \geq 2$, misalkan $X_j$
adalah banyaknya percobaan setelah keberhasilan ke-$(j - 1)$ hingga keberhasilan
ke-$j$.  Dapat ditunjukkan bahwa $X_1$,~$X_2$,~\dots merupakan proses percobaan
independen.
\begin{enumerate}
\item Apakah distribusi, nilai harapan, dan varians bersama bagi $X_j$?

\item Misalkan $T_n = X_1 + X_2 +\cdots+ X_n$.  Maka $T_n$ adalah waktu hingga
keberhasilan ke-$n$.  Carilah $E(T_n)$ dan $V(T_n)$.

\item Gunakan hasil (b) untuk mencari nilai harapan dan varians banyaknya lemparan
koin hingga kemunculan kepala ke-$n$.
\end{enumerate}

\i\label{exer 6.2.28} Dengan mengacu pada
Latihan~\ref{sec 6.1}.\ref{exer 6.1.30}, carilah varians banyaknya kotak Wheaties yang
dibeli sebelum memperoleh gambar setengah dari para pemain dan varians banyaknya kotak
tambahan yang diperlukan untuk memperoleh gambar separuh pemain sisanya.

\i\label{exer 6.2.99} Dalam Contoh~\ref{exam 5.3}, asumsikan bahwa buku yang dimaksud
memiliki 1000 halaman.  Misalkan $X$ adalah banyaknya halaman tanpa kesalahan.
Tunjukkan bahwa $E(X) = 905$ dan $V(X) = 86$.  Dengan menggunakan hasil ini, tunjukkan
bahwa peluang terdapat lebih dari 924 halaman tanpa kesalahan atau kurang dari 866
halaman tanpa kesalahan adalah ${} \leq .05$.

\i\label{exer 6.2.100} Misalkan $X$ berdistribusi Poisson dengan parameter
$\lambda$.  Tunjukkan bahwa $V(X) = \lambda$.

\end{LJSItem}

\choice{}{\section{Peubah Acak Kontinu}\label{sec 6.3}

Pada bagian ini kita membahas sifat-sifat nilai harapan dan varians suatu peubah acak
kontinu. Besaran-besaran ini didefinisikan sama seperti untuk peubah acak diskret dan
memiliki sifat-sifat yang sama.

\subsection*{Nilai Harapan}

\begin{definition}\label{def 6.5} Misalkan $X$ adalah peubah acak bernilai real dengan
fungsi densitas $f(x)$. \emx {Nilai harapan}\index{nilai harapan} $\mu = E(X)$
didefinisikan oleh
$$
\mu = E(X) = \int_{-\infty}^{+\infty} xf(x)\,dx\ ,
$$
asalkan integral
$$
\int_{-\infty}^{+\infty} |x|f(x)\,dx
$$
berhingga.
\end{definition}

Bandingkan definisi ini dengan definisi yang bersesuaian untuk peubah acak diskret
dalam Bagian~\ref{sec 6.1}. Secara intuitif, seperti pada bagian-bagian sebelumnya,
kita dapat menafsirkan $E(X)$ sebagai nilai yang kita harapkan akan diperoleh jika kita
melakukan banyak percobaan independen dan merata-ratakan nilai-nilai $X$ yang dihasilkan.

\par Sifat-sifat $E(X)$ dapat dirangkum sebagai berikut (bandingkan
Teorema~\ref{thm 6.1}).

\begin{theorem}\label{thm 6.10} Jika $X$ dan $Y$ adalah peubah acak bernilai real dan
$c$ adalah sebarang konstanta, maka
\begin{eqnarray*} 
E(X + Y) &=& E(X) + E(Y)\ , \\ 
E(cX)    &=& cE(X)\ .
\end{eqnarray*}
%\end{theorem}

Buktinya sangat mirip dengan bukti Teorema~\ref{thm 6.1}, sehingga tidak disertakan.
\end{theorem}
\par Secara lebih umum, jika $X_1$,~$X_2$, \dots,~$X_n$ adalah $n$ peubah acak
bernilai real dan $c_1$,~$c_2$, \dots,~$c_n$ adalah $n$ konstanta, maka
$$ E(c_1X_1 + c_2X_2 +\cdots+ c_nX_n) = c_1E(X_1) + c_2E(X_2) +\cdots+ c_nE(X_n)\ .
$$

\begin{example}\label{exam 6.16} Misalkan $X$ berdistribusi seragam pada interval
$[0, 1]$. Maka
$$E(X) = \int_0^1 x\,dx = 1/2\ .$$ Dengan demikian, jika kita memilih sejumlah besar
$N$ bilangan acak dari $[0,1]$ dan mengambil rataannya, kita dapat mengharapkan bahwa
rataan ini mendekati nilai harapan 1/2.
\end{example}

\begin{example}\label{exam 6.17} Misalkan $Z = (x, y)$ menyatakan sebuah titik yang
dipilih secara acak dan seragam dari cakram satuan, seperti dalam permainan lempar
panah pada Contoh~\ref{exam 2.2.2}, dan misalkan $X = (x^2 + y^2)^{1/2}$ adalah jarak
dari $Z$ ke pusat cakram. Dengan mudah dapat ditunjukkan bahwa fungsi densitas $X$
adalah $f(x) = 2x$, sehingga menurut definisi nilai harapan,
\begin{eqnarray*} E(X) & = & \int_0^1 x f(x)\,dx \\ & = & \int_0^1 x (2x)\,dx \\ & =
& \frac 23\ .
\end{eqnarray*}
\end{example}

\begin{example}\label{exam 6.18} Dalam contoh pasangan yang bertemu di penginapan
(Contoh~\ref{exam 2.2.7.4}), masing-masing orang tiba pada waktu yang berdistribusi
seragam antara pukul 5:00 dan 6:00 sore. Peubah acak $Z$ yang ditinjau adalah lamanya
orang pertama harus menunggu sampai orang kedua tiba. Telah ditunjukkan bahwa
$$f_Z(z) = 2(1-z)\ ,$$ untuk $0 \le z \le 1$. Jadi,
\begin{eqnarray*} E(Z) & = & \int_0^1 zf_Z(z)\,dz \\
     & = & \int_0^1 2z(1-z)\,dz \\
     & = & \Bigl[z^2 - \frac 23 z^3\Bigl]_0^1 \\  = \frac 13\ .
\end{eqnarray*}
\end{example}

\subsection*{Nilai Harapan Fungsi dari Peubah Acak}

Misalkan $X$ adalah peubah acak bernilai real dan $\phi(x)$ adalah fungsi kontinu dari
{\bf R} ke {\bf R}. Teorema berikut merupakan analog kontinu dari
Teorema~\ref{thm 6.3.5}.

\begin{theorem}\label{thm 6.11} Jika $X$ adalah peubah acak bernilai real dan jika
%$\phi : {\bf\rm R} \to {\bf\rm R}$ is a continuous real-valued function with domain
$\phi :$ {\bf R} $\to\ $ {\bf R} adalah fungsi kontinu bernilai real dengan domain
$[a,b]$, maka
$$ E(\phi(X)) = \int_{-\infty}^{+\infty} \phi(x) f_X(x)\, dx\ ,
$$ asalkan integral tersebut ada.
\end{theorem}

Untuk bukti teorema ini, lihat Ross.\index{ROSS, S.}\footnote{S. Ross,  \emx {A First Course in
Probability,} (New York: Macmillan, 1984), hlm. 241--245.}



\subsection*{Nilai Harapan Hasil Kali Dua Peubah Acak}

Secara umum, $E(XY) = E(X)E(Y)$ tidak selalu benar karena integral suatu hasil kali
bukanlah hasil kali integral-integralnya. Namun, jika $X$ dan $Y$ independen, nilai
harapan hasil kalinya sama dengan hasil kali nilai harapan masing-masing.

\begin{theorem}\label{thm 6.12} Misalkan $X$ dan $Y$ adalah peubah acak kontinu
bernilai real yang independen dengan nilai harapan berhingga. Maka
$$ E(XY) = E(X)E(Y)\ .
$$
\proof Kita hanya akan membuktikan kasus ketika daerah nilai $X$ dan $Y$ masing-masing
termuat dalam interval $[a, b]$ dan $[c, d]$. Misalkan fungsi densitas $X$ dan $Y$
masing-masing dinyatakan oleh $f_X(x)$ dan $f_Y(y)$. Karena $X$ dan $Y$ independen,
fungsi densitas bersama $X$ dan $Y$ adalah hasil kali fungsi-fungsi densitasnya. Jadi,
\begin{eqnarray*} E(XY) & = & \int_a^b \int_c^d xy f_X(x) f_Y(y)\, dy\,dx \\
      & = & \int_a^b x f_X(x)\, dx \int_c^d y f_Y(y)\, dy \\
      & = & E(X)E(Y)\ .
\end{eqnarray*}
\par
Bukti untuk kasus umum menggunakan barisan peubah acak terbatas yang mendekati $X$
dan $Y$ dan agak teknis, sehingga tidak disertakan.
\end{theorem}

Dengan cara yang sama, dapat ditunjukkan bahwa jika $X_1$,~$X_2$, \dots,~$X_n$
adalah $n$ peubah acak bernilai real yang saling independen, maka
$$ 
E(X_1 X_2 \cdots X_n) = E(X_1)\,E(X_2)\,\cdots\,E(X_n)\ .
$$

\begin{example}\label{exam 6.19} Misalkan $Z = (X, Y)$ adalah sebuah titik yang
dipilih secara acak dalam persegi satuan. Misalkan $A = X^2$ dan $B = Y^2$. Maka
Teorema~\ref{thm 4.3} mengakibatkan bahwa $A$ dan $B$ independen. Dengan menggunakan
Teorema~\ref{thm 6.11}, nilai harapan $A$ dan $B$ mudah dihitung:
\begin{eqnarray*} E(A) = E(B) & = & \int_0^1 x^2\,dx \\ & = & \frac 13\ .
\end{eqnarray*} Dengan menggunakan Teorema~\ref{thm 6.12}, nilai harapan $AB$ tidak
lain adalah hasil kali $E(A)$ dan $E(B)$, yaitu 1/9. Kegunaan teorema ini tampak jika
kita perhatikan bahwa menghitung $E(AB)$ dari definisi nilai harapan jauh lebih sulit.
Diperoleh bahwa fungsi densitas $AB$ adalah
$$f_{AB}(t) = \frac {-\log(t)}{4\sqrt t}\ ,$$ so
\begin{eqnarray*} E(AB) & = & \int_0^1 t f_{AB}(t)\,dt \\ & = & \frac 19\ .
\end{eqnarray*}
\end{example}

\begin{example}\label{exam 6.20} Sekali lagi, misalkan $Z = (X, Y)$ adalah sebuah
titik yang dipilih secara acak dalam persegi satuan dan misalkan $W = X + Y$. Maka
$Y$ dan $W$ tidak independen, dan kita mempunyai
\begin{eqnarray*} E(Y) & = &\frac 12\ , \\ E(W) & = & 1\ , \\ E(YW) & = & E(XY + Y^2)
= E(X)E(Y) + \frac 13 = \frac 7{12}
\ne E(Y)E(W)\ .
\end{eqnarray*}
\end{example}

Sekarang kita beralih ke varians.


\subsection*{Varians}

\begin{definition}\label{def 6.6} Misalkan $X$ adalah peubah acak bernilai real dengan
fungsi densitas $f(x)$. \emx {Varians}\index{varians} $\sigma^2 = V(X)$ didefinisikan oleh
$$
\sigma^2 = V(X) = E((X - \mu)^2)\ . 
$$
\end{definition}
Hasil berikut mudah diperoleh dari Teorema~\ref{thm 6.3.5}. Ada cara lain untuk
menghitung varians peubah acak kontinu yang biasanya sedikit lebih mudah. Cara
tersebut diberikan dalam Teorema~\ref{thm 6.14}.
\begin{theorem}\label{thm 6.13.5} Jika $X$ adalah peubah acak bernilai real dengan
$E(X) = \mu$, maka
$$
\sigma^2 = \int_{-\infty}^\infty (x - \mu)^2 f(x)\, dx\ .
$$
\end{theorem}
\par Semua sifat yang dinyatakan dalam tiga teorema berikut dibuktikan dengan cara
yang tepat sama dengan pembuktian teorema-teorema yang bersesuaian untuk peubah acak
diskret dalam Bagian~\ref{sec 6.2}.
\begin{theorem}\label{thm 6.13} Jika $X$ adalah peubah acak bernilai real yang
didefinisikan pada $\Omega$ dan $c$ adalah sebarang konstanta, maka (bandingkan
Teorema~\ref{thm 6.6})
\begin{eqnarray*}
V(cX)    &=& c^2 V(X)\ , \\ 
V(X + c) &=& V(X)\ .
\end{eqnarray*}
\end{theorem}

\begin{theorem}\label{thm 6.14} Jika $X$ adalah peubah acak bernilai real dengan
$E(X) = \mu$, maka (bandingkan Teorema~\ref{thm 6.7})
$$ V(X) = E(X^2) - \mu^2\ .
$$
\end{theorem}

\begin{theorem}\label{thm 6.15} Jika $X$ dan $Y$ adalah peubah acak bernilai real yang
independen pada $\Omega$, maka (bandingkan Teorema~\ref{thm 6.8})
$$ V(X + Y) = V(X) + V(Y)\ .
$$
\end{theorem}

\begin{example} (lanjutan Contoh~\ref{exam 6.16})\label{exam 6.18.5} Jika $X$
berdistribusi seragam pada $[0, 1]$, maka dengan menggunakan Teorema~\ref{thm 6.14}
kita memperoleh
$$  V(X) = \int_0^1 \Bigl(x - \frac 12 \Bigr)^2\, dx = \frac 1{12}\ .
$$
\end{example}

\begin{example}\label{exam 6.21} Misalkan $X$ adalah peubah acak berdistribusi
eksponensial dengan parameter $\lambda$. Maka fungsi densitas $X$ adalah
$$ f_X(x) = \lambda e^{-\lambda x}\ .
$$ Dari definisi nilai harapan dan integrasi parsial, kita memperoleh
\begin{eqnarray*} E(X) & = & \int_0^\infty x f_X(x)\, dx \\
     & = & \lambda \int_0^\infty x e^{-\lambda x}\, dx \\
     & = & \biggl.-xe^{-\lambda x}\biggr|_0^\infty + \int_0^\infty e^{-\lambda x}\,
dx \\
     & = & 0 + \biggl.\frac {e^{-\lambda x}}{-\lambda}\biggr|_0^\infty =
\frac 1\lambda\ .
\end{eqnarray*}

Dengan cara serupa, menggunakan Teorema~\ref{thm 6.11} dan \ref{thm 6.14}, kita
memperoleh
\begin{eqnarray*} V(X) & = & \int_0^\infty x^2 f_X(x)\, dx - \frac 1{\lambda^2} \\
     & = & \lambda \int_0^\infty x^2 e^{-\lambda x}\, dx - \frac 1{\lambda^2} \\
     & = & \biggl.-x^2 e^{-\lambda x}\biggr|_0^\infty + 2\int_0^\infty x e^{-\lambda
x}\, dx - \frac 1{\lambda^2} \\
     & = & \biggl.-x^2 e^{-\lambda x}\biggr|_0^\infty - \biggl.\frac {2xe^{-\lambda
x}}\lambda\biggr|_0^\infty - \biggl.\frac 2{\lambda^2} e^{-\lambda x}\biggr|_0^\infty
-
\frac 1{\lambda^2} =
\frac 2{\lambda^2} - \frac 1{\lambda^2} = \frac 1{\lambda^2}\ .
\end{eqnarray*} Dalam kasus ini, $E(X)$ dan $V(X)$ keduanya berhingga jika $\lambda > 0$.
\end{example}

\begin{example}\label{exam 6.22} Misalkan $Z$ adalah peubah acak normal baku dengan
fungsi densitas
$$ f_Z(x) = \frac 1{\sqrt{2\pi}} e^{-x^2/2}\ .
$$ Karena fungsi densitas ini simetris terhadap sumbu-$y$, mudah ditunjukkan bahwa
$$\int_{-\infty}^\infty x f_Z(x)\, dx$$ bernilai 0. Namun, perlu diingat bahwa nilai
harapan didefinisikan sebagai integral di atas hanya jika integral
$$\int_{-\infty}^\infty |x| f_Z(x)\, dx$$ berhingga. Integral ini sama dengan
$$2\int_0^\infty x f_Z(x)\, dx\ ,$$ yang dengan mudah dapat ditunjukkan berhingga.
Jadi, nilai harapan $Z$ adalah 0.
\par Untuk menghitung varians $Z$, kita mulai dengan menerapkan Teorema~\ref{thm 6.14}:
$$V(Z) =  \int_{-\infty}^{+\infty} x^2 f_Z(x)\, dx - \mu^2\ .$$ Jika kita menulis
$x^2$ sebagai $x\cdot x$ dan melakukan integrasi parsial, kita memperoleh
$$\biggl.\frac 1{\sqrt{2\pi}} (-x e^{-x^2/2})\biggr|_{-\infty}^{+\infty} + \frac
1{\sqrt{2\pi}} \int_{-\infty}^{+\infty} e^{-x^2/2}\, dx\ .
$$ Suku pertama di atas dapat ditunjukkan sama dengan 0 karena ketika $x \rightarrow
\pm \infty$, $e^{-x^2/2}$ mengecil lebih cepat daripada pertumbuhan $x$. Suku kedua
hanyalah densitas normal baku yang diintegralkan pada domainnya, sehingga nilainya
adalah 1. Oleh karena itu, varians densitas normal baku sama dengan 1.
\par Sekarang misalkan $X$ adalah peubah acak normal (tidak harus baku) dengan
parameter $\mu$ dan $\sigma$. Maka fungsi densitas $X$ adalah
$$ f_X(x) = \frac 1{\sqrt{2\pi}\sigma} e^{-(x - \mu)^2/2\sigma^2}\ .
$$ Kita dapat menulis $X = \sigma Z + \mu$, dengan $Z$ sebagai peubah acak normal
baku. Karena menurut perhitungan di atas $E(Z) = 0$ dan $V(Z) = 1$,
Teorema~\ref{thm 6.10} dan \ref{thm 6.13} mengakibatkan bahwa
\begin{eqnarray*} 
E(X)  &=& E(\sigma Z + \mu)\ =\ \mu\ , \\ 
V(X)  &=& V(\sigma Z + \mu)\ =\ \sigma^2\ .
\end{eqnarray*}
\end{example}

\begin{example}\label{exam 6.23} Misalkan $X$ adalah peubah acak kontinu dengan fungsi
densitas Cauchy
$$ f_X(x) = \frac {a}{\pi} \frac {1}{a^2 + x^2}\ .
$$ Maka nilai harapan $X$ tidak ada karena integral
$$
\frac a\pi \int_{-\infty}^{+\infty} \frac {|x|\,dx}{a^2 + x^2}
$$ divergen. Dengan demikian, varians $X$ juga tidak ada. Densitas yang variansnya
tidak terdefinisi, seperti densitas Cauchy, dalam sejumlah hal penting berperilaku
cukup berbeda dari densitas yang variansnya berhingga. Salah satu contoh perbedaan ini
akan kita lihat dalam Bagian~\ref{sec 8.2}.
\end{example}


\subsection*{Percobaan Independen}

\begin{corollary}\label{cor 6.1}                                                
Jika $X_1$,~$X_2$, \dots,~$X_n$ merupakan proses percobaan independen dari peubah
acak bernilai real, dengan $E(X_i) = \mu$ dan $V(X_i) = \sigma^2$, dan jika
\begin{eqnarray*} S_n & = & X_1 + X_2 +\cdots+ X_n\ , \\ A_n & = & \frac {S_n}n\ ,
\end{eqnarray*} 
maka
\begin{eqnarray*}
E(S_n) & = & n\mu\ ,\\
E(A_n) & = & \mu\ ,\\
V(S_n) & = & n\sigma^2\ ,\\
V(A_n) & = & \frac {\sigma^2} n\ .
\end{eqnarray*}
Akibatnya, jika kita menetapkan
$$ 
S_n^* = \frac {S_n - n\mu}{\sqrt{n\sigma^2}}\ ,
$$ 
maka
\begin{eqnarray*} 
E(S_n^*) & = & 0\ ,\\
V(S_n^*) & = & 1\ .
\end{eqnarray*}
Kita mengatakan bahwa $S_n^*$ adalah \emx {versi baku dari} $S_n$ (lihat
Latihan~\ref{exer 6.2.13} dalam Bagian~\ref{sec 6.2}).
\end{corollary}


\subsection*{Antrean}\index{antrean}

\begin{example}\label{exam 6.24} Mari kita tinjau kembali masalah antrean, yaitu
masalah pelanggan yang menunggu pelayanan dalam antrean (lihat Contoh~\ref{exam
5.21}). Sekali lagi kita menganggap bahwa pelanggan memasuki antrean sedemikian rupa
sehingga waktu antar-kedatangan merupakan peubah acak $X$ yang berdistribusi
eksponensial dengan fungsi densitas
$$ f_X(t) = \lambda e^{-\lambda t}\ .
$$ Maka nilai harapan waktu antar-kedatangan hanyalah $1/\lambda$ (lihat
Contoh~\ref{exam 6.21}), seperti yang dinyatakan dalam Contoh~\ref{exam 5.21}.
Kebalikannya, yaitu $\lambda$, sering disebut \emx {laju kedatangan.} \emx {Waktu
pelayanan} seseorang yang berada paling depan didefinisikan sebagai lamanya orang
tersebut berada di kepala antrean sebelum pergi. Kita menganggap pelanggan dilayani
sedemikian rupa sehingga waktu pelayanan merupakan peubah acak lain $Y$ yang
berdistribusi eksponensial dengan fungsi densitas
$$ f_X(t) = \mu e^{-\mu t}\ .
$$ Maka nilai harapan waktu pelayanan adalah
$$ E(X) = \int_0^\infty t f_X(t)\, dt = \frac 1\mu\ .
$$ Kebalikan nilai harapan ini, yaitu $\mu$, sering disebut {\em laju pelayanan.}
\par Berdasarkan pengalaman sehari-hari dengan antrean, kita menduga bahwa jika laju
pelayanan lebih besar daripada laju kedatangan, ukuran rata-rata antrean akan cenderung
stabil. Namun, jika laju pelayanan lebih kecil daripada laju kedatangan, panjang
antrean akan cenderung bertambah tanpa batas (lihat Gambar~\ref{fig 5.17}). Simulasi
dalam Contoh~\ref{exam 5.21} cenderung mendukung pengalaman sehari-hari tersebut. Kita
dapat menyatakan kesimpulan ini dengan lebih tepat jika kita mendefinisikan \emx
{intensitas lalu lintas} sebagai hasil kali
$$
\rho = ({\rm laju\ kedatangan})({\rm waktu\ pelayanan\ rata-rata}) = \frac \lambda\mu =
\frac {1/\mu}{1/\lambda}\ .
$$ Intensitas lalu lintas juga merupakan rasio waktu pelayanan rata-rata terhadap
waktu rata-rata antar-kedatangan. Jika intensitas lalu lintas kurang dari~1, antrean
akan berperilaku wajar, tetapi jika lebih dari~1, antrean akan membesar tanpa
batas. Dalam kasus kritis $\rho = 1$, dapat ditunjukkan bahwa antrean akan menjadi
panjang, tetapi selalu ada saat-saat ketika antrean kosong.\footnote{L.
Kleinrock,  \emx {Queueing Systems,} vol.~2 (New York: John Wiley and Sons, 1975).}

Jika intensitas lalu lintas kurang dari~1, kita dapat memandang panjang antrean sebagai
peubah acak $Z$ dengan nilai harapan berhingga,
$$ E(Z) = N\ .
$$

Waktu yang dihabiskan seorang pelanggan dalam antrean dapat dipandang sebagai peubah
acak $W$ dengan nilai harapan berhingga,
$$ E(W) = T\ .
$$ Kemudian dapat dikatakan bahwa ketika seorang pelanggan memasuki antrean, ia
mengharapkan menemukan $N$ orang di depannya, sedangkan ketika meninggalkan antrean,
ia mengharapkan menemukan $\lambda T$ orang di belakangnya. Karena dalam keadaan
setimbang kedua besaran ini seharusnya sama, kita mengharapkan bahwa
$$ N = \lambda T\ .
$$ Hubungan terakhir ini disebut \emx {hukum Little untuk antrean.}\index{hukum
Little untuk antrean}\footnote{ibid., hlm.~17.} Kita tidak akan membuktikannya di sini.
Sebuah bukti dapat ditemukan dalam Ross.\index{ROSS, S.}\footnote{S. M. Ross,  \emx
{Applied Probability Models with Optimization Applications,} (San Francisco:
Holden-Day, 1970)} Perhatikan bahwa dalam kasus ini kita menghitung waktu tunggu semua
pelanggan, termasuk mereka yang sama sekali tidak perlu menunggu. Dalam simulasi pada
Bagian~\ref{sec 4.2}, pelanggan-pelanggan ini tidak diperhitungkan.
\par Jika kita mengetahui nilai harapan panjang antrean, kita dapat menggunakan hukum
Little untuk memperoleh nilai harapan waktu tunggu karena
$$ T = \frac N\lambda\ .
$$ Panjang antrean adalah peubah acak dengan distribusi diskret. Distribusi ini dapat
diestimasi melalui simulasi dengan mencatat panjang antrean pada saat seorang pelanggan
tiba. Hasil simulasi ini (menggunakan program {\bf Queue}) ditampilkan dalam
Gambar~\ref{fig 6.5}.
\putfig{4truein}{PSfig6-5}{Distribusi panjang antrean.}{fig 6.5}
\par Perhatikan bahwa distribusi tersebut tampak sebagai distribusi geometrik. Dalam
kajian teori antrean ditunjukkan bahwa distribusi panjang antrean dalam keadaan setimbang
memang merupakan distribusi geometrik dengan
$$ s_j = (1 - \rho) \rho^j \qquad {\rm untuk}\  j = 0, 1, 2, \dots\ ,
$$ jika $\rho < 1$.
Nilai harapan peubah acak dengan distribusi ini adalah
$$ N = \frac \rho{(1 - \rho)}
$$ (lihat Contoh~\ref{exam 6.8}). Dengan demikian, menurut hasil Little, nilai harapan
waktu tunggu adalah
$$ T = \frac \rho{\lambda(1 - \rho)} = \frac 1{\mu - \lambda}\ ,
$$ dengan $\mu$ sebagai laju pelayanan, $\lambda$ sebagai laju kedatangan, dan $\rho$
sebagai intensitas lalu lintas.
\par 
Dalam simulasi kita, laju kedatangan adalah 1 dan laju pelayanan adalah 1.1. Jadi,
intensitas lalu lintas adalah $1/1.1 = 10/11$, nilai harapan ukuran antrean adalah
$$
\frac {10/11}{(1 - 10/11)} = 10\ ,
$$ dan nilai harapan waktu tunggu adalah
$$
\frac 1{1.1 - 1} = 10\ .
%*********The simulation values weren't very close to the theoretical ones.
%Do we want to rerun it?
$$ Dalam simulasi kita, ukuran rata-rata antrean adalah 8.19 dan waktu tunggu rata-rata
adalah 7.37. Dalam Gambar~\ref{fig 6.5.5}, ditampilkan histogram waktu tunggu.
Histogram ini menunjukkan bahwa densitas waktu tunggu bersifat eksponensial dengan
parameter $\mu - \lambda$, dan memang demikian halnya.
\putfig{4truein}{PSfig6-5-5}{Distribusi waktu tunggu dalam antrean.}{fig 6.5.5}
\end{example}

\exercises
\begin{LJSItem}

\i\label{exer 6.3.1} Misalkan $X$ adalah peubah acak dengan daerah nilai $[-1,1]$ dan
misalkan $f_X(x)$ adalah fungsi densitas $X$. Carilah $\mu(X)$ dan $\sigma^2(X)$ jika,
untuk $|x| < 1$,

\begin{enumerate}
\item $f_X(x) = 1/2$.

\item $f_X(x) = |x|$.

\item $f_X(x) = 1 - |x|$.

\item $f_X(x) = (3/2) x^2$.
\end{enumerate}


\i\label{exer 6.3.3} Misalkan $X$ adalah peubah acak dengan daerah nilai $[-1,1]$ dan
$f_X$ adalah fungsi densitasnya. Carilah $\mu(X)$ dan $\sigma^2(X)$ jika, untuk $|x| >
1$, $f_X(x) = 0$, dan untuk $|x| < 1$,

\begin{enumerate}
\item $f_X(x) = (3/4)(1 - x^2)$.

\item $f_X(x) = (\pi/4)\cos(\pi x/2)$.

\item $f_X(x) = (x + 1)/2$.

\item $f_X(x) = (3/8)(x + 1)^2$.
\end{enumerate}

\i\label{exer 6.3.5} Masa pakai bola lampu super ACME, diukur dalam jam, merupakan
peubah acak $T$ dengan fungsi densitas $f_T(t) = \lambda^2 t e^{-\lambda t}$, dengan
$\lambda = .05$. Berapakah nilai harapan masa pakai bola lampu ini? Berapakah
variansnya?

\i\label{exer 6.3.6} Misalkan $X$ adalah peubah acak dengan daerah nilai $[-1,1]$
dan fungsi densitas $f_X(x) = ax + b$ jika $|x| < 1$.

\begin{enumerate}
\item Tunjukkan bahwa jika $\int_{-1}^{+1} f_X(x)\, dx = 1$, maka $b = 1/2$.

\item Tunjukkan bahwa jika $f_X(x) \geq 0$, maka $-1/2 \leq a \leq 1/2$.

\item Tunjukkan bahwa $\mu = (2/3)a$, dan karenanya $-1/3 \leq \mu \leq 1/3$.

\item Tunjukkan bahwa $\sigma^2(X) = (2/3)b - (4/9)a^2 = 1/3 - (4/9)a^2$.
\end{enumerate}

\i\label{exer 6.3.7} Misalkan $X$ adalah peubah acak dengan daerah nilai $[-1,1]$
dan fungsi densitas $f_X(x) = ax^2 + bx + c$ jika $|x| < 1$ dan~0 selainnya.

\begin{enumerate}
\item Tunjukkan bahwa $2a/3 + 2c = 1$ (lihat Latihan~\ref{exer 6.3.6}).

\item Tunjukkan bahwa $2b/3 = \mu(X)$.

\item Tunjukkan bahwa $2a/5 + 2c/3 = \sigma^2(X)$.

\item Carilah $a$,~$b$, dan~$c$ jika $\mu(X) = 0$, $\sigma^2(X) = 1/15$, lalu
buatlah sketsa grafik $f_X$.

\item Carilah $a$,~$b$, dan~$c$ jika $\mu(X) = 0$, $\sigma^2(X) = 1/2$, lalu
buatlah sketsa grafik $f_X$.
\end{enumerate}

\i\label{exer 6.3.8} Misalkan $T$ adalah peubah acak dengan daerah nilai $[0,\infty]$
dan $f_T$ adalah fungsi densitasnya. Carilah $\mu(T)$ dan $\sigma^2(T)$ jika, untuk $t
< 0$, $f_T(t) = 0$, dan untuk $t > 0$,
\begin{enumerate}

\item $f_T(t) = 3e^{-3t}$.

\item $f_T(t) = 9te^{-3t}$.

\item $f_T(t) = 3/(1 + t)^4$.
\end{enumerate}

\i\label{exer 6.3.9} Misalkan $X$ adalah peubah acak dengan fungsi densitas $f_X$.
Dengan menggunakan kalkulus dasar, tunjukkan bahwa fungsi
$$
\phi(a) = E((X - a)^2)
$$ mencapai nilai minimumnya ketika $a = \mu(X)$, dan dalam kasus tersebut $\phi(a) =
\sigma^2(X)$.

\i\label{exer 6.3.10} Misalkan $X$ adalah peubah acak dengan rataan $\mu$ dan varians
$\sigma^2$. Misalkan $Y = aX^2 + bX + c$. Carilah nilai harapan $Y$.

\i\label{exer 6.3.11} Misalkan $X$,~$Y$, dan~$Z$ adalah peubah acak independen yang
masing-masing memiliki rataan $\mu$ dan varians $\sigma^2$.

\begin{enumerate}
\item Carilah nilai harapan dan varians $S = X + Y + Z$.

\item Carilah nilai harapan dan varians $A = (1/3)(X + Y + Z)$.

\item Carilah nilai harapan $S^2$ dan $A^2$.
\end{enumerate}

\i\label{exer 6.3.12} Misalkan $X$ dan $Y$ adalah peubah acak independen dengan
fungsi densitas seragam pada $[0,1]$. Carilah

\begin{enumerate}
\item $E(|X - Y|)$.

\item $E(\max(X,Y))$.

\item $E(\min(X,Y))$.

\item $E(X^2 + Y^2)$.

\item $E((X + Y)^2)$.
\end{enumerate}

\i\label{exer 6.3.13} Pilsdorff Beer Company\index{Pilsdorff Beer Company}
mengoperasikan armada truk di sepanjang jalan 100~mil dari Hangtown\index{Hangtown}
ke Dry Gulch\index{Dry Gulch}. Truk-truk tersebut sudah tua dan dapat mogok di
sebarang titik sepanjang jalan dengan peluang yang sama. Di manakah perusahaan harus
menempatkan garasi agar nilai harapan jarak dari lokasi mogok yang tipikal ke garasi
minimum? Dengan kata lain, jika $X$ adalah peubah acak yang menyatakan lokasi mogok,
diukur misalnya dari Hangtown, dan $b$ menyatakan lokasi garasi, pilihan $b$ manakah
yang meminimumkan $E(|X - b|)$? Sekarang anggaplah $X$ tidak berdistribusi seragam
pada $[0,100]$, melainkan memiliki fungsi densitas $f_X(x) = 2x/10{,}000$. Dalam
kasus ini, pilihan $b$ manakah yang meminimumkan $E(|X - b|)$?

\i\label{exer 6.3.14} Carilah $E(X^Y)$, dengan $X$ dan $Y$ sebagai peubah acak
independen yang seragam pada $[0, 1]$. Kemudian verifikasikan jawaban Anda melalui
simulasi.

\i\label{exer 6.3.15} Misalkan $X$ adalah peubah acak yang mengambil nilai-nilai
non-negatif dan memiliki fungsi distribusi $F(x)$. Tunjukkan bahwa
$$ E(X) = \int_0^\infty (1 - F(x))\, dx\ .
$$ 
 \emx {Petunjuk}: Lakukan integrasi parsial.

Ilustrasikan hasil ini dengan menghitung $E(X)$ melalui metode ini jika $X$ memiliki
distribusi eksponensial $F(x) = 1 - e^{-\lambda x}$ untuk $x \geq 0$, dan $F(x) = 0$
selainnya.

\i\label{exer 6.3.15.5} Misalkan $X$ adalah peubah acak kontinu dengan fungsi
densitas $f_X(x)$. Tunjukkan bahwa jika
$$
\int_{-\infty}^{+\infty} x^2 f_X(x)\, dx < \infty\ ,
$$ maka
$$
\int_{-\infty}^{+\infty} |x| f_X(x)\, dx < \infty\ .
$$ 
 \emx {Petunjuk}: Di luar interval $[-1, 1]$, integran pertama lebih besar daripada
integran kedua.

\i\label{exer 6.3.16} Misalkan $X$ adalah peubah acak yang berdistribusi seragam pada
$[0,20]$. Definisikan peubah acak baru $Y$ dengan $Y = \lfloor X\rfloor$ (bilangan
bulat terbesar yang tidak melebihi $X$). Carilah nilai harapan~$Y$. Lakukan hal yang sama untuk $Z =
\lfloor X + .5\rfloor$. Hitung $E\bigl(|X-Y|\bigr)$ dan $E\bigl(|X-Z|\bigr)$.
(Perhatikan bahwa $Y$ adalah nilai~$X$ yang dibulatkan ke bilangan bulat lebih kecil
terdekat, sedangkan $Z$ adalah nilai~$X$ yang dibulatkan ke bilangan bulat terdekat.
Metode pembulatan manakah yang lebih baik? Mengapa?)

\i\label{exer 6.3.17} Anggaplah masa pakai suatu komponen mesin diesel merupakan
peubah acak $X$ dengan densitas $f_X$. Ketika komponen tersebut aus, komponen itu
diganti dengan komponen lain yang memiliki densitas sama. Misalkan $N(t)$ adalah
banyaknya komponen yang digunakan selama waktu $t$. Kita ingin mempelajari peubah acak
$N(t)/t$. Karena rata-rata komponen diganti setiap $E(X)$ satuan waktu, kita
mengharapkan sekitar $t/E(X)$ komponen digunakan selama waktu $t$. Artinya, kita
mengharapkan bahwa
$$
\lim_{t \to \infty} E \Bigl(\frac {N(t)}t\Bigr) = \frac 1{E(X)}\ .
$$ Hasil ini benar, tetapi cukup sulit dibuktikan. Tulislah program yang memungkinkan
Anda menentukan densitas $f_X$ dan waktu $t$, lalu menyimulasikan percobaan ini untuk
memperoleh $N(t)/t$. Mintalah program Anda mengulangi percobaan tersebut 500 kali dan
menggambar diagram batang untuk hasil-hasil acak $N(t)/t$. Dari data ini, estimasikan
$E(N(t)/t)$ dan bandingkan dengan $1/E(X)$. Secara khusus, lakukan hal ini untuk~$t =
100$ dengan dua densitas berikut:

\begin{enumerate}
\item $f_X = e^{-t}$.

\item $f_X = te^{-t}$.
\end{enumerate}

\i\label{exer 6.3.18} Misalkan $X$ dan $Y$ adalah peubah acak. \emx {Kovarians}
$\rm {Cov}(X,Y)$ didefinisikan oleh (lihat Latihan~\ref{sec 6.2}.\ref{exer 6.2.24})
$$
\rm {cov}(X,Y) = E ((X - \mu(X))(Y - \mu(Y)) )\ .
$$
\begin{enumerate}
\item Tunjukkan bahwa $\rm {cov}(X,Y) = E(XY) - E(X)E(Y)$.

\item Dengan menggunakan (a), tunjukkan bahwa ${\rm cov}(X,Y) = 0$, jika $X$ dan
$Y$ independen. (Perhatian: kebalikannya \emx {tidak} selalu benar.)

\item Tunjukkan bahwa $V(X + Y) = V(X) + V(Y) + 2{\rm cov}(X,Y)$.
\end{enumerate}

\i\label{exer 6.3.19} Misalkan $X$ dan $Y$ adalah peubah acak dengan varians positif.
\emx {Korelasi} $X$ dan $Y$ didefinisikan sebagai
$$
\rho(X,Y) = \frac{{\rm cov}(X,Y)}{\sqrt{V(X)V(Y)}}\ .
$$
\begin{enumerate}
\item Dengan menggunakan Latihan~\ref{exer 6.3.18}(c), tunjukkan bahwa
$$ 0 \leq V\left( \frac X{\sigma(X)} + \frac Y{\sigma(Y)} \right) = 2(1 +
\rho(X,Y))\ .
$$

\item Sekarang tunjukkan bahwa
$$ 0 \leq V\left( \frac X{\sigma(X)} - \frac Y{\sigma(Y)} \right) = 2(1 -
\rho(X,Y))\ .
$$

\item Dengan menggunakan (a) dan (b), tunjukkan bahwa
$$ -1 \leq \rho(X,Y) \leq 1\ .
$$
\end{enumerate}

\i\label{exer 6.3.20} Misalkan $X$ dan $Y$ adalah peubah acak independen dengan
densitas seragam pada $[0,1]$. Misalkan $Z = X + Y$ dan $W = X - Y$. Carilah

\begin{enumerate}
\item $\rho(X,Y)$ (lihat Latihan~\ref{exer 6.3.19}).

\item $\rho(X,Z)$.

\item $\rho(Y,W)$.

\item $\rho(Z,W)$.
\end{enumerate}

\istar\label{exer 6.3.21} Ketika mempelajari data fisiologis tertentu, seperti tinggi
badan ayah dan anak laki-laki, sering kali wajar untuk menganggap bahwa data tersebut
(misalnya, tinggi badan para ayah dan anak laki-laki mereka) dideskripsikan oleh peubah
acak dengan densitas normal. Namun, peubah-peubah acak ini tidak independen, melainkan
berkorelasi. Sebagai contoh, densitas normal baku dua dimensi untuk peubah acak yang
berkorelasi berbentuk
$$ f_{X,Y}(x,y) = \frac 1{2\pi\sqrt{1 - \rho^2}} \cdot e^{-(x^2 - 2\rho xy + y^2)/2(1
- \rho^2)}\ .
$$
\begin{enumerate}
\item Tunjukkan bahwa $X$ dan $Y$ masing-masing memiliki densitas normal baku.

\item Tunjukkan bahwa korelasi $X$ dan $Y$ (lihat Latihan~\ref{exer 6.3.19}) adalah
$\rho$.
\end{enumerate}

\istar\label{exer 6.3.22} Untuk peubah acak berkorelasi $X$ dan $Y$, wajar untuk
menanyakan nilai harapan~$X$ jika~$Y$ diketahui. Sebagai contoh,
Galton\index{GALTON, F.} menghitung nilai harapan tinggi badan seorang anak laki-laki
jika tinggi badan ayahnya diketahui. Ia menggunakannya untuk menunjukkan bahwa pria
tinggi rata-rata dapat diharapkan memiliki anak laki-laki yang tidak setinggi mereka.
Demikian pula, siswa yang memperoleh hasil sangat baik pada suatu ujian dapat
diharapkan memperoleh hasil yang tidak sebaik itu pada ujian berikutnya, dan
seterusnya. Hal ini disebut \emx {regresi menuju rataan.}\index{regresi menuju rataan}
Untuk mendefinisikan nilai harapan bersyarat ini, pertama-tama kita mendefinisikan
densitas bersyarat~$X$ jika~$Y = y$ diketahui dengan
$$ f_{X|Y}(x|y) = \frac {f_{X,Y}(x,y)}{f_Y(y)}\ ,
$$ dengan $f_{X,Y}(x,y)$ sebagai densitas bersama $X$ dan $Y$, serta $f_Y$ sebagai
densitas $Y$. Kemudian nilai harapan bersyarat~$X$ jika~$Y$ diketahui adalah
$$ E(X|Y = y) = \int_a^b x f_{X|Y}(x|y)\, dx\ .
$$ Untuk densitas normal dalam Latihan~\ref{exer 6.3.21}, tunjukkan bahwa densitas
bersyarat $f_{X|Y}(x|y)$ adalah normal dengan rataan $\rho y$ dan varians $1 - \rho^2$.
Dari sini kita melihat bahwa jika $X$ dan $Y$ berkorelasi positif $(0 < \rho < 1)$ dan
jika $y > E(Y)$, maka nilai harapan~$X$ jika~$Y = y$ diketahui akan kurang dari~$y$
(yakni, terjadi regresi menuju rataan).

\i\label{exer 6.3.23} Sebuah titik $Y$ dipilih secara acak dari $[0,1]$. Kemudian
titik kedua $X$ dipilih dari interval $[0,Y]$. Carilah densitas $X$. \emx
{Petunjuk}: Hitung $f_{X|Y}$ seperti dalam Latihan~\ref{exer 6.3.22}, lalu gunakan
$$ f_X(x) = \int_x^1 f_{X|Y}(x|y) f_Y(y)\, dy\ .
$$ Dapatkah Anda juga menurunkan hasil tersebut secara geometris?

%**********Provide an answer in the answer key to the following problem  (it has been significantly
%changed).   In addition, we should find out what it means to be inside one of the ellipses
%discussed in the problem.

\istar\label{exer 6.3.24} Misalkan $X$ dan $V$ adalah dua peubah acak normal baku.
Misalkan $\rho$ adalah bilangan real antara -1 dan 1.
\begin{enumerate}
\item
Misalkan $Y = \rho X + \sqrt{1 - \rho^2} V$. Tunjukkan bahwa $E(Y) = 0$ dan $Var(Y)
= 1$. Kelak akan kita lihat (lihat Contoh~\ref{exam 7.8} dan
Contoh~\ref{exam 10.3.3}) bahwa jumlah dua peubah acak normal independen juga normal.
Jadi, dengan mengasumsikan fakta ini, kita telah menunjukkan bahwa $Y$ adalah normal
baku.
\item
Dengan menggunakan Latihan~\ref{exer 6.3.18} dan \ref{exer 6.3.19}, tunjukkan bahwa
korelasi $X$ dan $Y$ adalah $\rho$.
\item
Dalam Latihan~\ref{exer 6.3.21}, diberikan fungsi densitas bersama $f_{X,Y}(x, y)$
untuk peubah acak $(X, Y)$. Sekarang anggaplah kita ingin mengetahui himpunan titik
$(x, y)$ pada bidang-$xy$ sedemikian sehingga $f_{X,Y}(x, y) = C$ untuk suatu
konstanta $C$. Himpunan titik ini disebut himpunan berdensitas konstan. Secara kasar,
himpunan berdensitas konstan adalah himpunan titik tempat hasil $(X, Y)$ memiliki
peluang yang sama untuk jatuh. Tunjukkan bahwa untuk $C$ tertentu, himpunan titik
berdensitas konstan merupakan kurva dengan persamaan
$$x^2 - 2\rho x y + y^2 = D\ ,$$
dengan $D$ sebagai konstanta yang bergantung pada $C$. (Kurva ini merupakan elips.)
\item
Elips pada bagian (c) dapat digambar dengan menggunakan persamaan parametrik
\begin{eqnarray*} x & = & \frac {r\cos\theta}{\sqrt{2(1 - \rho)}} + \frac
{r\sin\theta}{\sqrt{2(1 +
\rho)}}\ , \\ y & = & \frac {r\cos\theta}{\sqrt{2(1 - \rho)}} - \frac
{r\sin\theta}{\sqrt{2(1 +
\rho)}}\ .
\end{eqnarray*}
Tulislah program untuk menggambar 1000 pasangan $(X, Y)$ bagi $\rho = -1/2, 0, 1/2$.
Pada setiap gambar, mintalah program Anda menggambar kurva parametrik di atas untuk
$r = 1, 2, 3$.
\end{enumerate}

\istar\label{exer 6.3.25} Mengikuti Galton, anggaplah bahwa tinggi badan para ayah dan
anak laki-laki mereka merupakan peubah acak normal yang dependen. Anggaplah tinggi
badan rata-rata adalah 68~inci, simpangan bakunya 2.7~inci, dan koefisien korelasinya
adalah~.5 (lihat Latihan~\ref{exer 6.3.21}~dan~\ref{exer 6.3.22}). Artinya, anggaplah
tinggi badan para ayah dan anak laki-laki mereka masing-masing berbentuk $2.7X + 68$
dan $2.7Y + 68$, dengan $X$ dan $Y$ sebagai peubah acak normal baku yang berkorelasi,
dengan koefisien korelasi~.5.

\begin{enumerate}
\item Berapakah nilai harapan tinggi badan anak laki-laki dari seorang ayah yang
tinggi badannya 72~inci?

\item Gambarlah diagram pencar tinggi badan 1000 pasangan ayah dan anak laki-laki.
\emx {Petunjuk}: Anda dapat memilih pasangan baku seperti dalam
Latihan~\ref{exer 6.3.24}, lalu menggambar $(2.7X + 68, 2.7Y + 68)$.
\end{enumerate}


%***********This problem is very time-consuming for the student to do.  Either we
%should make a file containing the data, or else we should delete the problem.  We
%will try to find the original data (instead of just the table).

\istar\label{exer 6.3.26} Jika kita memiliki pasangan data $(x_i,y_i)$ yang merupakan
hasil dari pasangan peubah acak dependen $X$,~$Y$, kita dapat mengestimasi koefisien
korelasi $\rho$ dengan
$$
\bar r = \frac {\sum_i (x_i - \bar x)(y_i - \bar y)}{(n - 1)s_Xs_Y}\ ,
$$ dengan $\bar x$ dan $\bar y$ masing-masing sebagai rataan sampel $X$ dan $Y$, serta
$s_X$ dan $s_Y$ sebagai simpangan baku sampel $X$ dan $Y$ (lihat
Latihan~\ref{sec 6.2}.\ref{exer 6.2.18}). Tulislah program untuk menghitung rataan
sampel, varians, dan korelasi bagi data dependen semacam itu. Gunakan program Anda
untuk menghitung besaran-besaran tersebut bagi data Galton tentang tinggi badan orang
tua dan anak yang diberikan dalam Lampiran~B.
\par
Gambarlah elips berdensitas sama sebagaimana didefinisikan dalam
Latihan~\ref{exer 6.3.24} untuk~$r = 4$,~6, dan~8, lalu pada grafik yang sama cetaklah
nilai-nilai yang muncul dalam tabel pada titik-titik yang sesuai. Sebagai contoh,
cetaklah 12 pada titik $(70.5,68.2)$, yang menunjukkan bahwa terdapat 12 kasus ketika
tinggi badan orang tua adalah 70.5 dan tinggi badan anak adalah 68.12. Periksalah
apakah data Galton konsisten dengan elips-elips berdensitas sama tersebut.

\i\label{exer 6.3.27} (dari Hamming\index{HAMMING, R. W.}\footnote{R.~W.~Hamming,
\emx {The Art of Probability for Scientists and Engineers} (Redwood City:
Addison-Wesley, 1991), hlm.~192.}) Anggaplah Anda berdiri di tepi sungai yang lurus.
\begin{enumerate}
\item Pilihlah secara acak suatu arah yang membuat Anda tetap berada di daratan, lalu
berjalanlah 1 km ke arah tersebut. Misalkan $P$ menyatakan posisi Anda. Berapakah nilai
harapan jarak dari $P$ ke sungai?
\item Sekarang anggaplah Anda bertindak seperti pada bagian (a), tetapi ketika tiba di
$P$, Anda memilih arah acak (dari \emx {semua} arah) dan berjalan sejauh 1 km.
Berapakah peluang Anda mencapai sungai sebelum perjalanan kedua selesai?
\end{enumerate}

\i\label{exer 6.3.28} (dari Hamming\index{HAMMING, R. W.}\footnote{ibid., hlm. 205.})
Sebuah permainan dimainkan sebagai berikut: Bilangan acak $X$ dipilih secara seragam
dari $[0, 1]$. Kemudian barisan bilangan acak $Y_1, Y_2, \ldots$ dipilih secara
independen dan seragam dari $[0, 1]$. Permainan berakhir untuk pertama kalinya ketika
$Y_i~>~X$. Anda kemudian dibayar $(i-1)$ dolar. Berapakah biaya masuk yang adil untuk
permainan ini?

\i\label{exer 6.3.29} Sebuah jarum panjang dengan panjang $L$ yang jauh lebih besar
daripada 1 dijatuhkan pada kisi yang garis horizontal dan vertikalnya berjarak satu
satuan. Tunjukkan bahwa banyaknya garis yang dipotong rata-rata, yaitu $a$, kira-kira
$$
a = \frac{4L}\pi\ .
$$

\end{LJSItem}}
%\end{document}
